%============================================================
%  Bd. 30 - Endliche Halbgruppen
%============================================================

\documentclass[main.tex]{subfiles}

\ifSubfilesClassLoaded{
  \BandSetupStandalone{B30}
}{
}

\title{Bd. 30 - Endliche Halbgruppen}
\author{Martin Kunze}
\date{}
\setcounter{file}{30}

\hypersetup{
  pdftitle={Bd. 30 - Endliche Halbgruppen},
  pdfauthor={Martin Kunze}
}

\begin{document}

\title{Bd. 30 - Endliche Halbgruppen}
\author{Martin Kunze}
\date{}
\hypersetup{pageanchor=false}
\maketitle
\hypersetup{pageanchor=true}
\setcounter{tocdepth}{3}
\tableofcontents
\setcounter{file}{30}
\renewcommand{\FormulaBandID}{30}

\chapter{Einführung}

Dieser Band untersucht Halbgruppen mit endlichem Träger und das
Isomorphieproblem ihrer Potenzhalbgruppen.  Nach den Grundbegriffen folgt
zunächst die Strukturtheorie: Der Index--Perioden-Satz folgt aus dem
Schubfachprinzip; daraus folgen idempotente
Potenzen und die Existenz von Idempotenten.  Die Idealtheorie aus Band~22
liefert anschließend das eindeutige kleinste nichtleere Ideal.  Erst danach
werden Multiplikationstafeln betrachtet.  Selbstbijektionen des Trägers
wirken durch Umbenennung auf den assoziativen Operationen, und ihre Orbits
sind genau die Isomorphieklassen.  Die Nummerierung der Potenzhalbgruppe wird
so in Blöcke zerlegt, dass jeder Eintrag unmittelbar auf ein Teilrechteck
der ursprünglichen Tafel zurückgeführt werden kann.  Ein weiteres Kapitel
fragt, welche Informationen sich aus einer Potenzhalbgruppe wieder über
ihre Ausgangshalbgruppe gewinnen lassen; die Schwierigkeit, Einermengen rein
multiplikativ zu erkennen, führt zur Potenzhalbgruppenvermutung.  Die drei
rekursiven Tafelverfahren stehen getrennt im rechnerischen Anhang.

\chapter{Endliche Halbgruppen}

\section{Definition}

\FormulaDefDeltaK[Begriff der endlichen Halbgruppe]{%
  \FiniteSemigroupAlg{A}{\star}%
}{FiniteSemigroupDef}{%
  \DeltaRow{Mengen}{A}
  \DeltaRow{\textbf{Axiome}}{}
  \DeltaPrem{Halbgruppe}{\SemigroupAlg{A}{\star}}
  \DeltaPrem{Endlichkeit}{\Finite{A}}
  \DeltaRow{\textbf{Neues Symbol}}{}
  \DeltaPrem{Endliche Halbgruppe}{\FiniteSemigroupAlg{A}{\star}}
}

\FormulaAxiomDeltaK[Zugriff auf die Halbgruppenstruktur]{%
  \FiniteSemigroupAlg{A}{\star}\vdash\SemigroupAlg{A}{\star}%
}{FiniteSemigroupBaseAxiom}{%
  \DeltaRow{Mengen}{A}
}

\FormulaAxiomDeltaK[Zugriff auf die Endlichkeit]{%
  \FiniteSemigroupAlg{A}{\star}\vdash\Finite{A}%
}{FiniteSemigroupFinitenessAxiom}{%
  \DeltaRow{Mengen}{A}
}

\section{Beispiele}

\FormulaThmDeltaK[Eine konstante Operation bildet auf einer Paarmenge eine
endliche Halbgruppe]{%
  \begin{aligned}[t]
    &\forall x,y\in\{p,q\}\,(x\star_0y=p)\vdash{}\\[-2pt]
    &\FiniteSemigroupAlg{\{p,q\}}{\star_0}
  \end{aligned}
}{ConstantFinitePairSemigroup}{%
  \DeltaRow{Mengen}{p\dsep q\dsep x\dsep y}
  \DeltaRow{Funktionen}{\star_0}
}
\begin{tabproof}
  \proofstep{1}{\forall x,y\in\{p,q\}\,(x\star_0y=p)}{\rA}
  \proofstep{}{p\in\{p,q\}}{\FormulaRefAuto{a\in\{a,b\}}}
  \proofstep{1}{\SemigroupWithZeroAlg{\{p,q\}}{\star_0}{p}}{%
    \FormulaRefAuto{ConstantSemigroupWithZero}[thm]{2,1}}
  \proofstep{1}{\SemigroupAlg{\{p,q\}}{\star_0}}{%
    \FormulaRefAuto{SemigroupWithZeroBaseAxiom}[ax]{3}}
  \proofstep{}{\Finite{\{p,q\}}}{\FormulaRefAuto{FinitePair}[thm]}
  \proofstep{1}{\FiniteSemigroupAlg{\{p,q\}}{\star_0}}{%
    \FormulaRefAuto{FiniteSemigroupDef}[def]{4,5}}
\end{tabproof}

Für $p\neq q$ besitzt der Träger dieses Beispiels genau zwei Elemente.

\section{Endlichkeit der Potenzhalbgruppe}

Der folgende Satz verbindet die zuvor eingeführte endliche Spezialisierung
mit dem Potenzhalbgruppenbeispiel.

\FormulaThmDeltaK[Potenzhalbgruppen endlicher Halbgruppen sind endlich]{%
  \FiniteSemigroupAlg{A}{\star}
  \vdash\FiniteSemigroupAlg{\PowNE{A}}{\star_{\mathcal P}}%
}{FinitePowerSemigroup}{%
  \DeltaRow{Mengen}{A}
  \DeltaRow{Funktionen}{\star\dsep\star_{\mathcal P}}
}
\begin{tabproof}
  \proofstep{1}{\FiniteSemigroupAlg{A}{\star}}{\rA}
  \proofstep{1}{\SemigroupAlg{A}{\star}}{%
    \FormulaRefAuto{FiniteSemigroupBaseAxiom}[ax]{1}}
  \proofstep{1}{\Finite{A}}{%
    \FormulaRefAuto{FiniteSemigroupFinitenessAxiom}[ax]{1}}
  \proofstep{1}{\SemigroupAlg{\PowNE{A}}{\star_{\mathcal P}}}{%
    \FormulaRefAuto{NonemptyPowerSemigroup}[thm]{2}}
  \proofstep{}{\PowNE{A}\subseteq\powerset(A)}{%
    \FormulaRefAuto{A\setminus B\subseteq A}[thm]}
  \proofstep{1}{\Finite{\powerset(A)}}{%
    \FormulaRefAuto{FinitePowerSet}[thm]{3}}
  \proofstep{1}{\Finite{\PowNE{A}}}{%
    \FormulaRefAuto{FiniteSubset}[thm]{5,6}}
  \proofstep{1}{\FiniteSemigroupAlg{\PowNE{A}}{\star_{\mathcal P}}}{%
    \FormulaRefAuto{FiniteSemigroupDef}[def]{4,7}}
\end{tabproof}

Die genaue Anzahl der nichtleeren Teilmengen halten wir in der im ersten
Band eingeführten Tabellenform fest.

\FormulaThmDeltaK[Anzahl der nichtleeren Teilmengen]{%
  \Finite{A}\dsep N\in\mathbb N\dsep \tCard(A)=N
  \vdash \tCard\bigl(\PowNE{A}\bigr)=2^N-1%
}{FiniteSemigroupNonemptyPowerSetCardinality}{%
  \DeltaRow{Mengen}{A\dsep N}
}
\begin{tabproofwide}
  \proofstepwidestar[1]{\Finite{A}}{\rA}
  \proofstepwidestar[2]{N\in\mathbb N}{\rA}
  \proofstepwidestar[3]{\tCard(A)=N}{\rA}
  \proofstepwidestarbreak[1,2,3]{%
    \tCard\bigl(\PowNE{A}\bigr)+1=2^N\land
    \tCard\bigl(\PowNE{A}\bigr)=2^N-1}{%
    \FormulaRefAuto{FiniteNonemptyPowerSetCardinality}[thm]{1,2,3}}
  \proofstepwidestarforcebreak[1,2,3]{%
    \tCard\bigl(\PowNE{A}\bigr)=2^N-1}{\rAEb{4}}
\end{tabproofwide}

Für \(\tCard(A)=N\) gilt also
\(\left|\powerset(A)\setminus\{\varnothing\}\right|=2^N-1\).
Aus einer dreielementigen Grundhalbgruppe entsteht insbesondere eine
Potenzhalbgruppe mit \(2^3-1=7\) Elementen.  Diese Größenrechnung erklärt
bereits die sieben Zeilen und Spalten der späteren Potenzhalbgruppentafel.

\section{Morphismen}

\FormulaDefDeltaK[Homomorphismus endlicher Halbgruppen]{%
  \FiniteSemigroupHom{f}{A,\star}{B,\diamond}%
}{FiniteSemigroupHomDef}{%
  \DeltaRow{Mengen}{A\dsep B}
  \DeltaRow{Funktionen}{f\dsep\star\dsep\diamond}
  \DeltaRow{\textbf{Axiome}}{}
  \DeltaPrem{Quellstruktur}{\FiniteSemigroupAlg{A}{\star}}
  \DeltaPrem{Zielstruktur}{\FiniteSemigroupAlg{B}{\diamond}}
  \DeltaPrem{Halbgruppenhomomorphismus}{%
    \SemigroupHom{f}{A,\star}{B,\diamond}}
  \DeltaRow{\textbf{Neues Symbol}}{}
  \DeltaPrem{Homomorphismus}{%
    \FiniteSemigroupHom{f}{A,\star}{B,\diamond}}
}

\FormulaAxiomDeltaK[Quellstruktur]{%
  \FiniteSemigroupHom{f}{A,\star}{B,\diamond}
  \vdash\FiniteSemigroupAlg{A}{\star}%
}{FiniteSemigroupHomSourceAxiom}{%
  \DeltaRow{Mengen}{A\dsep B}
  \DeltaRow{Funktionen}{f\dsep\star\dsep\diamond}
}

\FormulaAxiomDeltaK[Zielstruktur]{%
  \FiniteSemigroupHom{f}{A,\star}{B,\diamond}
  \vdash\FiniteSemigroupAlg{B}{\diamond}%
}{FiniteSemigroupHomTargetAxiom}{%
  \DeltaRow{Mengen}{A\dsep B}
  \DeltaRow{Funktionen}{f\dsep\star\dsep\diamond}
}

\FormulaAxiomDeltaK[Zugriff auf den Halbgruppenhomomorphismus]{%
  \FiniteSemigroupHom{f}{A,\star}{B,\diamond}
  \vdash\SemigroupHom{f}{A,\star}{B,\diamond}%
}{FiniteSemigroupHomBaseHomAxiom}{%
  \DeltaRow{Mengen}{A\dsep B}
  \DeltaRow{Funktionen}{f\dsep\star\dsep\diamond}
}

\FormulaDefDeltaK[Isomorphismus endlicher Halbgruppen]{%
  \FiniteSemigroupIso{f}{A,\star}{B,\diamond}%
}{FiniteSemigroupIsoDef}{%
  \DeltaRow{Mengen}{A\dsep B}
  \DeltaRow{Funktionen}{f\dsep\star\dsep\diamond}
  \DeltaRow{\textbf{Axiome}}{}
  \DeltaPrem{Homomorphismus}{%
    \FiniteSemigroupHom{f}{A,\star}{B,\diamond}}
  \DeltaPrem{Bijektivität}{f\colon A\bij B}
  \DeltaRow{\textbf{Neues Symbol}}{}
  \DeltaPrem{Isomorphismus}{%
    \FiniteSemigroupIso{f}{A,\star}{B,\diamond}}
}

\FormulaAxiomDeltaK[Homomorphismuszugriff]{%
  \FiniteSemigroupIso{f}{A,\star}{B,\diamond}
  \vdash\FiniteSemigroupHom{f}{A,\star}{B,\diamond}%
}{FiniteSemigroupIsoHomAxiom}{%
  \DeltaRow{Mengen}{A\dsep B}
  \DeltaRow{Funktionen}{f\dsep\star\dsep\diamond}
}

\FormulaAxiomDeltaK[Bijektivitätszugriff]{%
  \FiniteSemigroupIso{f}{A,\star}{B,\diamond}
  \vdash f\colon A\bij B%
}{FiniteSemigroupIsoBijectiveAxiom}{%
  \DeltaRow{Mengen}{A\dsep B}
  \DeltaRow{Funktionen}{f\dsep\star\dsep\diamond}
}

%============================================================
\chapter{Struktur endlicher Halbgruppen}

\ProofWideReasonsNearFormula
%============================================================

Die allgemeinen Begriffe der positiven Potenz, der erzeugten
Unterhalbgruppe, des idempotenten Elements und des Ideals stehen bereits in
Band~22 zur Verfügung.  Die Endlichkeit wird nun nur noch an den Stellen
verwendet, an denen sie tatsächlich neue Struktur erzwingt: Eine unendliche
Potenzfolge mit endlichem Wertebereich muss sich wiederholen, und die
nichtleeren Ideale bilden eine nichtleere Familie von Teilmengen eines
endlichen Trägers.

\section{Index und Periode}

\subsection{Wiederholungsstellen}

Die Potenzfolge aus Band~22 beginnt beim Folgenindex \(0\), während positive
Potenzen beim Exponenten \(1\) beginnen.  Die Wiederholungsstellen werden
deshalb von Anfang an mit positiven Exponenten nummeriert.  Nur im folgenden
Existenzbeweis wird ein nullbasierter Treffer \(k<n\) einmalig zu den
positiven Exponenten \(k+1<n+1\) verschoben; danach ist keine Umrechnung mehr
nötig.

\FormulaDefDeltaK[Wiederholungsstellen einer Potenzfolge]{%
  \begin{aligned}[t]
  &\SemigroupAlg{A}{\star}\dsep a\in A\vdash{}\\[-2pt]
  &\mathcal R_{A,\star}(a)\coloneqq
  \bigl\{j\in\mathbb N\bigm|{}
    j\neq0\land
    \exists i\in\mathbb N\,
      \bigl(i\neq0\land i<j\land a^i=a^j\bigr)
  \bigr\}
  \end{aligned}%
}{FiniteSemigroupPowerRepetitionSetDef}{%
  \DeltaRow{Mengen}{A\dsep a\dsep i\dsep j}
  \DeltaRow{Funktionen}{\star}
  \DeltaRow{\textbf{Neues Symbol}}{}
  \DeltaPrem{Wiederholungsstellen}{\mathcal R_{A,\star}(a)}
}

\FormulaThmDeltaK[Die Potenzfolge eines Elements einer endlichen
Halbgruppe wiederholt sich]{%
  \FiniteSemigroupAlg{A}{\star}\dsep a\in A
  \vdash \mathcal R_{A,\star}(a)\neq\varnothing%
}{FiniteSemigroupPowerRepetitionExists}{%
  \DeltaRow{Mengen}{A\dsep a\dsep i\dsep j\dsep k\dsep n}
  \DeltaRow{Funktionen}{\star}
}
\begin{tabproofwide}
  \proofstepwidestar[1]{\FiniteSemigroupAlg{A}{\star}}{\rA}
  \proofstepwidestar[2]{a\in A}{\rA}
  \proofstepwidestar[1]{\SemigroupAlg{A}{\star}}{%
    \FormulaRefAuto{FiniteSemigroupBaseAxiom}[ax]{1}}
  \proofstepwidestar[1]{\Finite{A}}{%
    \FormulaRefAuto{FiniteSemigroupFinitenessAxiom}[ax]{1}}
  \proofstepwidestarforcebreak[1,2]{%
    \bigl(a^{n+1}\bigr)_{n\in\mathbb N}\colon\mathbb N\to A}{%
    \FormulaRefAuto{SemigroupPowerSequenceFunction}[thm]{3,2},
    \FormulaRefAuto{SemigroupPowerSuccessorIndex}[thm]{3,2}}
  \proofstepwidestarbreak[1,2]{%
    \exists k\,\exists n\,
      \bigl(k\in\mathbb N\land n\in\mathbb N\land k<n\land
        a^{k+1}=a^{n+1}\bigr)}{%
    \FormulaRefAuto{FiniteSequenceRepetition}[thm]{4,5}}
  \proofstepwidestar[7]{%
    \exists n\,\bigl(k\in\mathbb N\land n\in\mathbb N\land
      k<n\land a^{k+1}=a^{n+1}\bigr)}{\rA}
  \proofstepwidestar[8]{%
    k\in\mathbb N\land n\in\mathbb N\land k<n\land
      a^{k+1}=a^{n+1}}{\rA}
  \proofstepwidestar[8]{1\in\mathbb N}{%
    \FormulaRefAuto{PeanoOneInN}[thm]}
  \proofstepwidestar[8]{k\in\mathbb N}{\rAEa{8}}
  \proofstepwidestar[8]{n\in\mathbb N}{\rAEa{\rAEb{8}}}
  \proofstepwidestarbreak[8]{k+1\in\mathbb N}{%
    \FormulaRefAuto{PeanoAddClosure}[thm]{10,9}}
  \proofstepwidestarbreak[8]{n+1\in\mathbb N}{%
    \FormulaRefAuto{PeanoAddClosure}[thm]{11,9}}
  \proofstepwidestarbreak[8]{k+1\neq0}{%
    \rIE{%
      \FormulaRefAuto{PeanoAddOneSucc}[thm]{10},
      \FormulaRefAuto{PeanoSuccNonzero}[ax]{10}}}
  \proofstepwidestarbreak[8]{n+1\neq0}{%
    \rIE{%
      \FormulaRefAuto{PeanoAddOneSucc}[thm]{11},
      \FormulaRefAuto{PeanoSuccNonzero}[ax]{11}}}
  \proofstepwidestarbreak[8]{1+k<1+n}{%
    \FormulaRefAuto{PeanoAddLeftStrictMonotone}[thm]{%
      9,10,11,\rAEa{\rAEb{\rAEb{8}}}}}
  \proofstepwidestarbreak[8]{k+1=1+k}{%
    \FormulaRefAuto{PeanoAddCommutative}[thm]{10,9}}
  \proofstepwidestarbreak[8]{n+1=1+n}{%
    \FormulaRefAuto{PeanoAddCommutative}[thm]{11,9}}
  \proofstepwidestarbreak[8]{k+1<n+1}{\rIE{17,18,16}}
  \proofstepwidestarbreak[8]{%
    \exists i\in\mathbb N\,
      \bigl(i\neq0\land i<n+1\land a^i=a^{n+1}\bigr)}{%
    \rEI{\rAI{%
      12,\rAI{14,\rAI{19,\rAEb{\rAEb{\rAEb{8}}}}}}}}
  \proofstepwidestarbreak[8]{%
    \exists j\in\mathbb N\,
      \Bigl(j\neq0\land
        \exists i\in\mathbb N\,
          \bigl(i\neq0\land i<j\land a^i=a^j\bigr)\Bigr)}{%
    \rEI{\rAI{13,\rAI{15,20}}}}
  \proofstepwidestarbreak[7]{%
    \exists j\in\mathbb N\,
      \Bigl(j\neq0\land
        \exists i\in\mathbb N\,
          \bigl(i\neq0\land i<j\land a^i=a^j\bigr)\Bigr)}{%
    \rEE{7,8,21}}
  \proofstepwidestarbreak[1,2]{%
    \exists j\in\mathbb N\,
      \Bigl(j\neq0\land
        \exists i\in\mathbb N\,
          \bigl(i\neq0\land i<j\land a^i=a^j\bigr)\Bigr)}{%
    \rEE{6,7,22}}
  \proofstepwidestar[24]{%
    j\in\mathbb N\land
      \Bigl(j\neq0\land
        \exists i\in\mathbb N\,
          \bigl(i\neq0\land i<j\land a^i=a^j\bigr)\Bigr)}{\rA}
  \proofstepwidestar[24]{%
    \exists i\in\mathbb N\,
      \bigl(i\neq0\land i<j\land a^i=a^j\bigr)}{%
    \rAEb{\rAEb{24}}}
  \proofstepwidestar[26]{%
    i\in\mathbb N\land
      \bigl(i\neq0\land i<j\land a^i=a^j\bigr)}{\rA}
  \proofstepwidestarbreak[1,2,24,26]{j\in\mathcal R_{A,\star}(a)}{%
    \FormulaRefAuto{FiniteSemigroupPowerRepetitionSetDef}[def]{%
      3,2,\rAEa{24},\rAEa{\rAEb{24}},\rEI{26}}}
  \proofstepwidestar[1,2,24,26]{%
    \mathcal R_{A,\star}(a)\neq\varnothing}{%
    \FormulaRefAuto{x\in A\vdash A\neq\varnothing}[thm]{27}}
  \proofstepwidestar[1,2,24]{\mathcal R_{A,\star}(a)\neq\varnothing}{%
    \rEE{25,26,28}}
  \proofstepwidestar[1,2]{\mathcal R_{A,\star}(a)\neq\varnothing}{%
    \rEE{23,24,29}}
\end{tabproofwide}

\subsection{Axiomatischer Aufbau des Index--Perioden-Paares}

Die Bezeichnung \emph{Index--Perioden-Paar} ist in der Halbgruppentheorie
üblich und wird deshalb beibehalten.  Anschaulich ist \(\mu\) der
\emph{Zykluseintritt} und \(\lambda\) die \emph{Zykluslänge}: Die erste
Wiederholung endet bei \(\mu+\lambda\), und dort kehrt der Potenzwert an der
Stelle \(\mu\) wieder.  Wir führen den Begriff axiomatisch ein; insbesondere
sind die drei benötigten Eigenschaften anschließend unmittelbar zugänglich.

\FormulaAxiomDeltaK[Axiom der positiven natürlichen Parameter]{%
  \begin{aligned}[t]
    &\SemigroupAlg{A}{\star}\dsep a\in A\dsep
      \operatorname{IP}_{A,\star}(a;\mu,\lambda)\vdash{}\\[-2pt]
    &\mu\in\mathbb N\land\mu\neq0\land
      \lambda\in\mathbb N\land\lambda\neq0
  \end{aligned}%
}{FiniteSemigroupIndexPeriodPositiveData}{%
  \DeltaRow{Mengen}{A\dsep a\dsep\mu\dsep\lambda}
  \DeltaRow{Funktionen}{\star}
}

\FormulaAxiomDeltaK[Axiom der Wiederholungsgleichung]{%
  \SemigroupAlg{A}{\star}\dsep a\in A\dsep
  \operatorname{IP}_{A,\star}(a;\mu,\lambda)
  \vdash a^{\mu+\lambda}=a^\mu%
}{FiniteSemigroupIndexPeriodEquality}{%
  \DeltaRow{Mengen}{A\dsep a\dsep\mu\dsep\lambda}
  \DeltaRow{Funktionen}{\star}
}

\FormulaAxiomDeltaK[Trennungsaxiom vor dem ersten Wiederholungsendpunkt]{%
  \SemigroupAlg{A}{\star}\dsep a\in A\dsep
  \operatorname{IP}_{A,\star}(a;\mu,\lambda)\dsep
  r,s\in\mathbb N\dsep r\neq0\dsep r<s\dsep s<\mu+\lambda
  \vdash a^r\neq a^s%
}{FiniteSemigroupIndexPeriodInitialDistinctness}{%
  \DeltaRow{Mengen}{A\dsep a\dsep\mu\dsep\lambda\dsep r\dsep s}
  \DeltaRow{Funktionen}{\star}
  \DeltaRow{Relationsoperatoren}{<_{\mathbb N}}
}

\FormulaDefDeltaK[Axiomatischer Begriff des Index--Perioden-Paares]{%
  \operatorname{IP}_{A,\star}(a;\mu,\lambda)%
}{FiniteSemigroupIndexPeriodPairDef}{%
  \DeltaRow{Mengen}{A\dsep a\dsep\mu\dsep\lambda\dsep r\dsep s}
  \DeltaRow{Funktionen}{\star}
  \DeltaRow{Relationsoperatoren}{<_{\mathbb N}}
  \DeltaRow{\textbf{Axiome}}{}
  \DeltaRow{Index}{\mu\in\mathbb N\land\mu\neq0}
  \DeltaRow{Periode}{\lambda\in\mathbb N\land\lambda\neq0}
  \DeltaRow{Wiederholung am Endpunkt}{a^{\mu+\lambda}=a^\mu}
  \DeltaRow{Trennung}{%
    \begin{aligned}[t]
      &r,s\in\mathbb N\dsep r\neq0\dsep r<s\dsep s<\mu+\lambda\\[-2pt]
      &\Longrightarrow a^r\neq a^s
    \end{aligned}}
  \DeltaRow{\textbf{Neues Symbol}}{}
  \DeltaPrem{Index--Perioden-Paar}{%
    \operatorname{IP}_{A,\star}(a;\mu,\lambda)}
}

Da \(\mathcal R_{A,\star}(a)\) nun die positiven Wiederholungsendpunkte
enthält, bedeutet \enquote{vor der ersten Wiederholung} genau \(s<j\).
Die Schranke \(s<j+1\) würde den bereits wiederholten Endpunkt \(j\) selbst
einschließen.

\FormulaThmDeltaK[Minimalität einer Wiederholungsstelle trennt alle
vorherigen positiven Potenzen]{%
  \begin{aligned}[t]
    &\SemigroupAlg{A}{\star}\dsep a\in A\dsep j\in\mathbb N\dsep
      \forall v\in\mathcal R_{A,\star}(a)\,(j\leq v)\dsep{}\\[-2pt]
    &r,s\in\mathbb N\dsep r\neq0\dsep r<s\dsep s<j
      \vdash a^r\neq a^s
  \end{aligned}%
}{FiniteSemigroupLeastRepetitionInitialDistinctness}{%
  \DeltaRow{Mengen}{A\dsep a\dsep j\dsep r\dsep s}
  \DeltaRow{Funktionen}{\star}
  \DeltaRow{Relationsoperatoren}{<_{\mathbb N}\dsep\leq_{\mathbb N}}
}
\begin{tabproofwide}
  \proofstepwidestar[1]{\SemigroupAlg{A}{\star}}{\rA}
  \proofstepwidestar[2]{a\in A}{\rA}
  \proofstepwidestar[3]{j\in\mathbb N}{\rA}
  \proofstepwidestar[4]{%
    \forall v\in\mathcal R_{A,\star}(a)\,(j\leq v)}{\rA}
  \proofstepwidestar[5]{r,s\in\mathbb N}{\rA}
  \proofstepwidestar[6]{r\neq0}{\rA}
  \proofstepwidestar[7]{r<s}{\rA}
  \proofstepwidestar[8]{s<j}{\rA}
  \proofstepwidestar[9]{a^r=a^s}{\rA}
  \proofstepwidestarbreak[5,7]{s\neq0}{%
    \FormulaRefAuto{PeanoStrictUpperNonzero}[thm]{%
      \rAEa{5},\rAEb{5},7}}
  \proofstepwidestarbreak[1,2,5--7,9]{%
    s\in\mathcal R_{A,\star}(a)}{%
    \FormulaRefAuto{FiniteSemigroupPowerRepetitionSetDef}[def]{%
      1,2,\rAEb{5},10,
      \rEI{\rAI{\rAEa{5},\rAI{6,\rAI{7,9}}}}}}
  \proofstepwidestarbreak[1,2,4--7,9]{j\leq s}{%
    \rRE{11,\rUE{4}}}
  \proofstepwidestarbreak[1--9]{\bot}{%
    \FormulaRefAuto{PeanoOrderNoCycle}[thm]{\rAEb{5},3,8,12}}
  \proofstepwidestar[1--8]{a^r\neq a^s}{\rCI{9,13}}
\end{tabproofwide}

\FormulaThmDeltaK[Die erste Wiederholungsstelle erzeugt ein
Index--Perioden-Paar]{%
  \begin{aligned}[t]
    &\SemigroupAlg{A}{\star}\dsep a\in A\dsep i,j,q\in\mathbb N\dsep
      i\neq0\dsep j\neq0\dsep i<j\dsep a^i=a^j\dsep{}\\[-2pt]
    &\forall v\in\mathcal R_{A,\star}(a)\,(j\leq v)\dsep
      i+(q+1)=j\vdash
      \operatorname{IP}_{A,\star}(a;i,q+1)
  \end{aligned}%
}{FiniteSemigroupLeastRepetitionProducesIndexPeriod}{%
  \DeltaRow{Mengen}{A\dsep a\dsep i\dsep j\dsep q\dsep r\dsep s}
  \DeltaRow{Funktionen}{\star}
  \DeltaRow{Relationsoperatoren}{<_{\mathbb N}\dsep\leq_{\mathbb N}}
}
\begin{tabproofwide}
  \proofstepwidestar[1]{\SemigroupAlg{A}{\star}}{\rA}
  \proofstepwidestar[2]{a\in A}{\rA}
  \proofstepwidestar[3]{i,j,q\in\mathbb N}{\rA}
  \proofstepwidestar[4]{i\neq0}{\rA}
  \proofstepwidestar[5]{j\neq0}{\rA}
  \proofstepwidestar[6]{i<j}{\rA}
  \proofstepwidestar[7]{a^i=a^j}{\rA}
  \proofstepwidestar[8]{%
    \forall v\in\mathcal R_{A,\star}(a)\,(j\leq v)}{\rA}
  \proofstepwidestar[9]{i+(q+1)=j}{\rA}
  \proofstepwidestarbreak[3]{q+1\in\mathbb N}{%
    \FormulaRefAuto{PeanoAddClosure}[thm]{%
      \rAEb{\rAEb{3}},\FormulaRefAuto{PeanoOneInN}[thm]}}
  \proofstepwidestarbreak[3]{q+1\neq0}{%
    \rIE{%
      \FormulaRefAuto{PeanoAddOneSucc}[thm]{\rAEb{\rAEb{3}}},
      \FormulaRefAuto{PeanoSuccNonzero}[ax]{\rAEb{\rAEb{3}}}}}
  \proofstepwidestarbreak[1,2,3,7,9]{a^{i+(q+1)}=a^i}{%
    \rChain{%
      \rIE{9},
      \FormulaRefAuto{a=b\vdash b=a}[thm]{7}}}
  \proofstepwidestar[13]{r\in\mathbb N}{\rA}
  \proofstepwidestar[14]{s\in\mathbb N}{\rA}
  \proofstepwidestar[15]{%
    r\neq0\land r<s\land s<i+(q+1)}{\rA}
  \proofstepwidestarbreak[3,9,15]{s<j}{%
    \rIE{9,\rAEb{\rAEb{15}}}}
  \proofstepwidestarbreak[1,2,3,8,13,14,15]{a^r\neq a^s}{%
    \FormulaRefAuto{FiniteSemigroupLeastRepetitionInitialDistinctness}[thm]{%
      1,2,\rAEa{\rAEb{3}},8,\rAI{13,14},
      \rAEa{15},\rAEa{\rAEb{15}},16}}
  \proofstepwidestarbreak[1,2,3,8,9]{%
    \forall r,s\in\mathbb N\,
      \bigl(r\neq0\land r<s\land s<i+(q+1)
        \rightarrow a^r\neq a^s\bigr)}{%
    \rUI{\rRI{13,\rUI{\rRI{14,\rRI{15,17}}}}}}
  \proofstepwidestarbreak[1,2,3,4,7,8,9]{%
    \operatorname{IP}_{A,\star}(a;i,q+1)}{%
    \FormulaRefAuto{FiniteSemigroupIndexPeriodPairDef}[def]{%
      1,2,\rAEa{3},4,10,11,12,18}}
\end{tabproofwide}

\FormulaThmDeltaK[Existenz eines Index--Perioden-Paares]{%
  \FiniteSemigroupAlg{A}{\star}\dsep a\in A
  \vdash
  \exists\mu\,\exists\lambda\,
    \operatorname{IP}_{A,\star}(a;\mu,\lambda)%
}{FiniteSemigroupIndexPeriodExistence}{%
  \DeltaRow{Mengen}{A\dsep a\dsep i\dsep j\dsep q\dsep
    \mu\dsep\lambda\dsep r\dsep s}
  \DeltaRow{Funktionen}{\star}
  \DeltaRow{Relationsoperatoren}{<_{\mathbb N}\dsep\leq_{\mathbb N}}
}

\begin{tabproofwide}
  \proofstepwidestar[1]{\FiniteSemigroupAlg{A}{\star}}{\rA}
  \proofstepwidestar[2]{a\in A}{\rA}
  \proofstepwidestar[1]{\SemigroupAlg{A}{\star}}{%
    \FormulaRefAuto{FiniteSemigroupBaseAxiom}[ax]{1}}
  \proofstepwidestar[1,2]{\mathcal R_{A,\star}(a)\neq\varnothing}{%
    \FormulaRefAuto{FiniteSemigroupPowerRepetitionExists}[thm]{1,2}}
  \proofstepwidestar[1,2]{\mathcal R_{A,\star}(a)\subseteq\mathbb N}{%
    \FormulaRefAuto{FiniteSemigroupPowerRepetitionSetDef}[def]{3,2}}
  \proofstepwidestar[1,2]{%
    \exists j\in\mathcal R_{A,\star}(a)\,
      \forall v\in\mathcal R_{A,\star}(a)\,(j\leq v)}{%
    \FormulaRefAuto{NaturalWellOrderingPrinciple}[thm]{5,4}}
  \proofstepwidestar[7]{%
    j\in\mathcal R_{A,\star}(a)\land
      \forall v\in\mathcal R_{A,\star}(a)\,(j\leq v)}{\rA}
  \proofstepwidestar[1,2,7]{j\in\mathbb N}{%
    \FormulaRefAuto{A\subseteq B\dsep x\in A\vdash x\in B}[thm]{%
      5,\rAEa{7}}}
  \proofstepwidestarbreak[1,2,7]{j\neq0}{%
    \FormulaRefAuto{FiniteSemigroupPowerRepetitionSetDef}[def]{%
      3,2,\rAEa{7}}}
  \proofstepwidestarbreak[1,2,7]{%
    \exists i\in\mathbb N\,
      \bigl(i\neq0\land i<j\land a^i=a^j\bigr)}{%
    \FormulaRefAuto{FiniteSemigroupPowerRepetitionSetDef}[def]{%
      3,2,\rAEa{7}}}
  \proofstepwidestar[11]{%
    i\in\mathbb N\land
      \bigl(i\neq0\land i<j\land a^i=a^j\bigr)}{\rA}
  \proofstepwidestarbreak[1,2,7,11]{%
    \exists q\in\mathbb N\,\bigl(i+(q+1)=j\bigr)}{%
    \FormulaRefAuto{PeanoStrictAddOneDistance}[thm]{%
      \rAEa{11},8,\rAEa{\rAEb{\rAEb{11}}}}}
  \proofstepwidestar[13]{q\in\mathbb N\land i+(q+1)=j}{\rA}
  \proofstepwidestarbreak[1,2,7,11,13]{%
    i,j,q\in\mathbb N}{%
    \rAI{\rAEa{11},\rAI{8,\rAEa{13}}}}
  \proofstepwidestar[1,2,7,11,13]{i\neq0}{\rAEa{\rAEb{11}}}
  \proofstepwidestar[1,2,7,11,13]{j\neq0}{9}
  \proofstepwidestar[1,2,7,11,13]{i<j}{%
    \rAEa{\rAEb{\rAEb{11}}}}
  \proofstepwidestarbreak[1,2,7,11,13]{a^i=a^j}{%
    \rAEb{\rAEb{\rAEb{11}}}}
  \proofstepwidestarbreak[1,2,7,11,13]{%
    \forall v\in\mathcal R_{A,\star}(a)\,(j\leq v)}{\rAEb{7}}
  \proofstepwidestarbreak[1,2,7,11,13]{i+(q+1)=j}{\rAEb{13}}
  \proofstepwidestarbreak[1,2,7,11,13]{%
    \operatorname{IP}_{A,\star}(a;i,q+1)}{%
    \FormulaRefAuto{FiniteSemigroupLeastRepetitionProducesIndexPeriod}[thm]{%
      3,2,14,15,16,17,18,19,20}}
  \proofstepwidestar[1,2,7,11,13]{%
    \exists\mu\,\exists\lambda\,
      \operatorname{IP}_{A,\star}(a;\mu,\lambda)}{\rEI{\rEI{21}}}
  \proofstepwidestar[1,2,7,11]{%
    \exists\mu\,\exists\lambda\,
      \operatorname{IP}_{A,\star}(a;\mu,\lambda)}{\rEE{12,13,22}}
  \proofstepwidestar[1,2,7]{%
    \exists\mu\,\exists\lambda\,
      \operatorname{IP}_{A,\star}(a;\mu,\lambda)}{\rEE{10,11,23}}
  \proofstepwidestar[1,2]{%
    \exists\mu\,\exists\lambda\,
      \operatorname{IP}_{A,\star}(a;\mu,\lambda)}{\rEE{6,7,24}}
\end{tabproofwide}

\FormulaThmDeltaK[Der Endpunkt eines Index--Perioden-Paares ist eindeutig]{%
  \begin{aligned}[t]
    &\SemigroupAlg{A}{\star}\dsep a\in A\dsep
      \operatorname{IP}_{A,\star}(a;\mu,\lambda)\dsep
      \operatorname{IP}_{A,\star}(a;\mu',\lambda')\vdash{}\\[-2pt]
    &\mu+\lambda=\mu'+\lambda'
  \end{aligned}%
}{FiniteSemigroupIndexPeriodEndpointUnique}{%
  \DeltaRow{Mengen}{A\dsep a\dsep\mu\dsep\lambda\dsep
    \mu'\dsep\lambda'}
  \DeltaRow{Funktionen}{\star}
  \DeltaRow{Relationsoperatoren}{<_{\mathbb N}\dsep\leq_{\mathbb N}}
}
\begin{proof}
Setze \(E=\mu+\lambda\) und \(E'=\mu'+\lambda'\).  Wäre
\(E<E'\), so lägen die gleichen Potenzen \(a^\mu=a^E\) vollständig vor
dem Endpunkt \(E'\): Es gilt \(\mu\neq0\), \(\mu<E<E'\).  Dies widerspricht
dem Trennungsaxiom des zweiten Index--Perioden-Paares.  Gälte dagegen
\(E'<E\), so lägen die gleichen Potenzen
\(a^{\mu'}=a^{E'}\) mit \(\mu'\neq0\) und
\(\mu'<E'<E\) vollständig vor dem Endpunkt \(E\); dies widerspräche dem
Trennungsaxiom des ersten Index--Perioden-Paares.  Die Trichotomie der
natürlichen Ordnung liefert daher \(E=E'\).
\end{proof}

\FormulaThmDeltaK[Der Index bei festem Wiederholungsendpunkt ist eindeutig]{%
  \begin{aligned}[t]
    &\SemigroupAlg{A}{\star}\dsep a\in A\dsep
      \operatorname{IP}_{A,\star}(a;\mu,\lambda)\dsep
      \operatorname{IP}_{A,\star}(a;\mu',\lambda')\dsep
      \mu+\lambda=\mu'+\lambda'\vdash{}\\[-2pt]
    &\mu=\mu'
  \end{aligned}%
}{FiniteSemigroupIndexPeriodIndexUnique}{%
  \DeltaRow{Mengen}{A\dsep a\dsep\mu\dsep\lambda\dsep
    \mu'\dsep\lambda'}
  \DeltaRow{Funktionen}{\star}
  \DeltaRow{Relationsoperatoren}{<_{\mathbb N}}
}
\begin{proof}
Schreibe wieder \(E=\mu+\lambda=\mu'+\lambda'\).  Aus den beiden
Wiederholungsgleichungsaxiomen folgt \(a^\mu=a^E=a^{\mu'}\).  Außerdem liegen
\(\mu\) und \(\mu'\) wegen der positiven Perioden beide strikt vor \(E\).
Wäre \(\mu<\mu'\) oder \(\mu'<\mu\), so widerspräche die erhaltene
Potenzgleichheit dem Trennungsaxiom vor \(E\).  Nach der Trichotomie ist
also \(\mu=\mu'\).
\end{proof}

\FormulaThmDeltaK[Eindeutigkeit des Index--Perioden-Paares]{%
  \begin{aligned}[t]
    &\SemigroupAlg{A}{\star}\dsep a\in A\dsep
      \operatorname{IP}_{A,\star}(a;\mu,\lambda)\dsep
      \operatorname{IP}_{A,\star}(a;\mu',\lambda')\vdash{}\\[-2pt]
    &\mu=\mu'\land\lambda=\lambda'
  \end{aligned}%
}{FiniteSemigroupIndexPeriodUniqueness}{%
  \DeltaRow{Mengen}{A\dsep a\dsep\mu\dsep\lambda\dsep
    \mu'\dsep\lambda'}
  \DeltaRow{Funktionen}{\star}
  \DeltaRow{Relationsoperatoren}{<_{\mathbb N}}
}

\begin{proof}
Der Endpunktsatz liefert
\(\mu+\lambda=\mu'+\lambda'\), und der anschließende Eindeutigkeitssatz
liefert \(\mu=\mu'\).  Durch Kürzen dieser gleichen linken Summanden folgt
\(\lambda=\lambda'\).
\end{proof}

\FormulaThmDeltaK[Index--Perioden-Satz für endliche Halbgruppen]{%
  \begin{aligned}[t]
    &\FiniteSemigroupAlg{A}{\star}\dsep a\in A\vdash{}\\[-2pt]
    &\exists! z\,\exists\mu\,\exists\lambda\,
      \bigl(z=(\mu,\lambda)\land
        \operatorname{IP}_{A,\star}(a;\mu,\lambda)\bigr)
  \end{aligned}%
}{FiniteSemigroupIndexPeriod}{%
  \DeltaRow{Mengen}{A\dsep a\dsep z\dsep z'\dsep
    \mu\dsep\lambda\dsep\mu'\dsep\lambda'}
  \DeltaRow{Funktionen}{\star}
}
\begin{proof}
Wähle nach dem Existenzsatz ein Index--Perioden-Paar
\((\mu,\lambda)\) und setze \(z=(\mu,\lambda)\).  Seien ferner
\(z=(\mu,\lambda)\) und \(z'=(\mu',\lambda')\) durch
Index--Perioden-Paare dargestellt.  Deren Eindeutigkeit liefert
\(\mu=\mu'\) und \(\lambda=\lambda'\), also \(z=z'\).
\end{proof}

Die eindeutig bestimmten Zahlen \(\mu\) und \(\lambda\) heißen
\emph{Index} und \emph{Periode} von \(a\).

\section{Periodische Fortsetzung und idempotente Potenzen}

Aus \(a^{\mu+\lambda}=a^\mu\) folgt mittels des Additionsgesetzes für
positive Potenzen die Gleichheit
\(a^{\mu+\lambda+t}=a^{\mu+t}\) für jedes \(t\in\mathbb N\).

\FormulaThmDeltaK[Fortsetzung einer Potenzgleichheit]{%
  \begin{aligned}[t]
    &\SemigroupAlg{A}{\star}\dsep a\in A\dsep
      \mu,\lambda\in\mathbb N\dsep\mu\neq0\dsep\lambda\neq0\dsep
      a^{\mu+\lambda}=a^\mu\dsep t\in\mathbb N\vdash{}\\[-2pt]
    &a^{(\mu+\lambda)+t}=a^{\mu+t}
  \end{aligned}%
}{FiniteSemigroupPowerEqualityPropagation}{%
  \DeltaRow{Mengen}{A\dsep a\dsep\mu\dsep\lambda\dsep t}
  \DeltaRow{Funktionen}{\star}
}
\begin{proof}
Für \(t=0\) ist die Behauptung genau die vorausgesetzte Potenzgleichheit.
Für \(t\neq0\) gilt nach dem Additionsgesetz
\[
  a^{(\mu+\lambda)+t}
  =a^{\mu+\lambda}\star a^t
  =a^\mu\star a^t
  =a^{\mu+t}.
\]
Damit sind beide Fälle erfasst.
\end{proof}

\FormulaThmDeltaK[Verschiebung um eine Periode]{%
  \begin{aligned}[t]
    &\SemigroupAlg{A}{\star}\dsep a\in A\dsep
      \mu,\lambda,r\in\mathbb N\dsep
      \mu\neq0\dsep\lambda\neq0\dsep
      a^{\mu+\lambda}=a^\mu\dsep\mu\leq r\vdash{}\\[-2pt]
    &a^{r+\lambda}=a^r
  \end{aligned}%
}{FiniteSemigroupSinglePeriodShift}{%
  \DeltaRow{Mengen}{A\dsep a\dsep\mu\dsep\lambda\dsep r\dsep t}
  \DeltaRow{Funktionen}{\star}
  \DeltaRow{Relationsoperatoren}{\leq_{\mathbb N}}
}
\begin{proof}
Aus \(\mu\leq r\) folgt \(r=\mu+t\) für ein \(t\in\mathbb N\).  Wegen
der Assoziativität und Kommutativität der Addition gilt
\((\mu+t)+\lambda=(\mu+\lambda)+t\).  Die Fortsetzung der
Potenzgleichheit ergibt daher
\[
  a^{r+\lambda}
  =a^{(\mu+t)+\lambda}
  =a^{(\mu+\lambda)+t}
  =a^{\mu+t}
  =a^r.
\]
\end{proof}

Sei \((\mu,\lambda)\) ein Index--Perioden-Paar.  Dann gilt
\(a^{\mu+\lambda}=a^\mu\), und jede Verschiebung des Exponenten ab \(\mu\)
um \(\lambda\) lässt den Potenzwert unverändert.

\FormulaThmDeltaK[Periodische Fortsetzung eines Index--Perioden-Paars]{%
  \begin{aligned}[t]
    &\SemigroupAlg{A}{\star}\dsep a\in A\dsep
      \operatorname{IP}_{A,\star}(a;\mu,\lambda)\vdash{}\\[-2pt]
    &\forall r\in\mathbb N\,
      \bigl(\mu\leq r\rightarrow a^{r+\lambda}=a^r\bigr)
  \end{aligned}%
}{FiniteSemigroupPowerEventuallyPeriodic}{%
  \DeltaRow{Mengen}{A\dsep a\dsep\mu\dsep\lambda\dsep r}
  \DeltaRow{Funktionen}{\star}
  \DeltaRow{Relationsoperatoren}{\leq_{\mathbb N}}
}
\begin{proof}
Aus der Definition des Index--Perioden-Paares erhält man die positiven
natürlichen Zahlen \(\mu,\lambda\) und die Gleichheit
\(a^{\mu+\lambda}=a^\mu\).  Für jedes \(r\geq\mu\) ist deshalb der Satz über
die Verschiebung um eine Periode anwendbar.
\end{proof}

\FormulaThmDeltaK[Größe der von einem Element erzeugten Unterhalbgruppe]{%
  \begin{aligned}[t]
    &\SemigroupAlg{A}{\star}\dsep a\in A\dsep
      \operatorname{IP}_{A,\star}(a;\mu,\lambda)\vdash{}\\[-2pt]
    &\langle a\rangle_{A,\star}
      =\bigl\{a^r\bigm|r\in\mathbb N\land r\neq0\land
        r<\mu+\lambda\bigr\}\land{}\\[-2pt]
    &\Finite{\langle a\rangle_{A,\star}}\land
      \tCard\!\left(\langle a\rangle_{A,\star}\right)
        =\mu+\lambda-1
  \end{aligned}%
}{FiniteSemigroupMonogenicSize}{%
  \DeltaRow{Mengen}{A\dsep a\dsep\mu\dsep\lambda\dsep r}
  \DeltaRow{Funktionen}{\star}
  \DeltaRow{Relationsoperatoren}{<_{\mathbb N}}
}
\begin{proof}
Nach Band~22 besteht \(\langle a\rangle_{A,\star}\) aus allen positiven
Potenzen von \(a\).  Wir zeigen zunächst, dass zu jedem positiven Exponenten
\(r\) ein positives \(s<\mu+\lambda\) mit \(a^r=a^s\) existiert.  Gäbe es
ein Gegenbeispiel, so hätte die Menge seiner Exponenten nach dem
Wohlordnungsprinzip ein kleinstes Element \(r\).  Dieses erfüllt
\(\mu+\lambda\leq r\), denn kleinere Exponenten repräsentieren sich selbst.
Die additive Charakterisierung der natürlichen Ordnung liefert ein
\(t\in\mathbb N\) mit
\[
  r=(\mu+\lambda)+t=(\mu+t)+\lambda.
\]
Für \(s_0=\mu+t\) gilt \(\mu\leq s_0<r\).  Die Verschiebung um eine
Periode ergibt \(a^r=a^{s_0}\).  Wegen der Minimalität von \(r\) lässt sich
\(a^{s_0}\) aber bereits durch einen positiven Exponenten
\(s<\mu+\lambda\) darstellen, ein Widerspruch.  Somit werden
alle positiven Potenzen bereits von den Exponenten
\(0<r<\mu+\lambda\) erfasst.  Diese Potenzen sind nach der
Paarverschiedenheitsbedingung paarweise verschieden.  Die beschränkte
  Exponentenmenge ist endlich und besitzt genau \(\mu+\lambda-1\)
Elemente.  Ihre Potenzabbildung ist auf ihr injektiv und hat die erzeugte
Unterhalbgruppe als Bild.  Damit ist diese Unterhalbgruppe endlich und ihre
  endliche Kardinalzahl gleich \(\mu+\lambda-1\).
\end{proof}

\FormulaThmDeltaK[Verschiebung um ein Vielfaches der Periode]{%
  \begin{aligned}[t]
    &\SemigroupAlg{A}{\star}\dsep a\in A\dsep
      \mu,\lambda,r,q\in\mathbb N\dsep
      \mu\neq0\dsep\lambda\neq0\dsep
      a^{\mu+\lambda}=a^\mu\dsep\mu\leq r\vdash{}\\[-2pt]
    &a^{r+\lambda\cdot q}=a^r
  \end{aligned}%
}{FiniteSemigroupMultiplePeriodShift}{%
  \DeltaRow{Mengen}{A\dsep a\dsep\mu\dsep\lambda\dsep r\dsep q\dsep k}
  \DeltaRow{Funktionen}{\star}
  \DeltaRow{Relationsoperatoren}{\leq_{\mathbb N}}
}

\begin{proof}
Wir induzieren über \(q\).  Für \(q=0\) ist
\(r+\lambda\cdot0=r\).  Gelte die Behauptung für \(k\).  Wegen
\(\mu\leq r\leq r+\lambda\cdot k\) darf die Verschiebung um eine Periode am
Exponenten \(r+\lambda\cdot k\) angewandt werden.  Damit folgt
\[
\begin{aligned}
a^{r+\lambda\cdot(k+1)}
 &=a^{(r+\lambda\cdot k)+\lambda}\\
 &=a^{r+\lambda\cdot k}
 =a^r.
\end{aligned}
\]
Das Induktionsprinzip liefert die Behauptung für alle \(q\in\mathbb N\).
\end{proof}

\FormulaThmDeltaK[Ein geeignetes Periodenvielfaches ist positiv und liegt
ab dem Index]{%
  \begin{aligned}[t]
    &\mu,\lambda\in\mathbb N\dsep\mu\neq0\dsep\lambda\neq0\vdash{}\\[-2pt]
    &\lambda\cdot\mu\in\mathbb N\land
      \lambda\cdot\mu\neq0\land
      \mu\leq\lambda\cdot\mu
  \end{aligned}%
}{FiniteSemigroupIdempotentExponentData}{%
  \DeltaRow{Mengen}{\mu\dsep\lambda}
  \DeltaRow{Relationsoperatoren}{\leq_{\mathbb N}}
}
\begin{proof}
Die Abgeschlossenheit der Multiplikation liefert
\(\lambda\cdot\mu\in\mathbb N\).  Aus der Nullteilerfreiheit und
\(\lambda,\mu\neq0\) folgt außerdem
\(\lambda\cdot\mu\neq0\).  Schreibe wegen \(\lambda\neq0\)
\(\lambda=k+1\).  Dann ist
\[
  \lambda\cdot\mu
  =(k+1)\cdot\mu
  =k\cdot\mu+\mu
  =\mu+k\cdot\mu,
\]
also \(\mu\leq\lambda\cdot\mu\).
\end{proof}

\FormulaThmDeltaK[Jedes Element einer endlichen Halbgruppe besitzt eine
idempotente Potenz]{%
  \begin{aligned}[t]
    &\FiniteSemigroupAlg{A}{\star}\dsep a\in A\vdash{}\\[-2pt]
    &\exists n\in\mathbb N\,
      \bigl(n\neq0\land a^n\in E_{A,\star}\bigr)
  \end{aligned}%
}{FiniteSemigroupIdempotentPower}{%
  \DeltaRow{Mengen}{A\dsep a\dsep\mu\dsep\lambda\dsep n}
  \DeltaRow{Funktionen}{\star}
}

\begin{proof}
Wähle das Index--Perioden-Paar \((\mu,\lambda)\) von \(a\) und setze
\(n=\lambda\cdot\mu\).  Der vorherige Satz liefert
\(n\in\mathbb N\), \(n\neq0\) und \(\mu\leq n\).  Die Verschiebung um
\(\mu\) Perioden, angewandt am Exponenten \(n\), ergibt
\[
  a^{n+n}
  =a^{n+\lambda\cdot\mu}
  =a^n.
\]
Nach dem Additionsgesetz ist \(a^{n+n}=a^n\star a^n\).  Da positive
Potenzen in \(A\) liegen, ist \(a^n\) somit ein idempotentes Element.
\end{proof}

\FormulaThmDeltaK[Existenz eines Idempotenten in einer nichtleeren
endlichen Halbgruppe]{%
  \FiniteSemigroupAlg{A}{\star}\dsep A\neq\varnothing
  \vdash E_{A,\star}\neq\varnothing%
}{FiniteSemigroupIdempotentExists}{%
  \DeltaRow{Mengen}{A\dsep a\dsep n}
  \DeltaRow{Funktionen}{\star}
}
\begin{tabproofwide}
  \proofstepwidestar[1]{\FiniteSemigroupAlg{A}{\star}}{\rA}
  \proofstepwidestar[2]{A\neq\varnothing}{\rA}
  \proofstepwidestar[2]{\exists a\in A}{%
    \FormulaRefAuto{A\neq\varnothing\vdash\exists x\in A}[thm]{2}}
  \proofstepwidestar[4]{a\in A}{\rA}
  \proofstepwidestarbreak[1,4]{%
    \exists n\in\mathbb N\,
      \bigl(n\neq0\land a^n\in E_{A,\star}\bigr)}{%
    \FormulaRefAuto{FiniteSemigroupIdempotentPower}[thm]{1,4}}
  \proofstepwidestar[6]{%
    n\in\mathbb N\land n\neq0\land a^n\in E_{A,\star}}{\rA}
  \proofstepwidestar[6]{E_{A,\star}\neq\varnothing}{%
    \FormulaRefAuto{x\in A\vdash A\neq\varnothing}[thm]{%
      \rAEb{\rAEb{6}}}}
  \proofstepwidestar[1,4]{E_{A,\star}\neq\varnothing}{\rEE{5,6,7}}
  \proofstepwidestar[1,2]{E_{A,\star}\neq\varnothing}{\rEE{3,4,8}}
\end{tabproofwide}

\section{Das kleinste nichtleere Ideal}

\FormulaDefDeltaK[Familie der nichtleeren zweiseitigen Ideale]{%
  \SemigroupAlg{A}{\star}\vdash
  \mathcal I_{A,\star}\coloneqq
  \bigl\{I\in\powerset(A)\bigm|\mathsf{Id}(I;A,\star)\bigr\}%
}{FiniteSemigroupIdealFamilyDef}{%
  \DeltaRow{Mengen}{A\dsep I}
  \DeltaRow{Funktionen}{\star}
  \DeltaRow{\textbf{Neues Symbol}}{}
  \DeltaPrem{Idealfamilie}{\mathcal I_{A,\star}}
}

\FormulaThmDeltaK[Existenz eines kleinsten nichtleeren Ideals]{%
  \begin{aligned}[t]
    &\FiniteSemigroupAlg{A}{\star}\dsep A\neq\varnothing\vdash{}\\[-2pt]
    &\exists K\,
      \Bigl(\mathsf{Id}(K;A,\star)\land
        \forall I\,\bigl(\mathsf{Id}(I;A,\star)\rightarrow K\subseteq I\bigr)
      \Bigr)
  \end{aligned}%
}{FiniteSemigroupLeastIdealExistence}{%
  \DeltaRow{Mengen}{A\dsep I\dsep J\dsep K}
  \DeltaRow{Funktionen}{\star}
}
\begin{proof}
Nach \FormulaRefAuto{SemigroupCarrierIsIdeal}[thm] gehört \(A\) wegen
\(A\neq\varnothing\) zur Idealfamilie.  Diese Familie ist eine Teilmenge von
\(\powerset(A)\); ihre Endlichkeit folgt daher aus
\FormulaRefAuto{FinitePowerSet}[thm] und
\FormulaRefAuto{FiniteSubset}[thm]. Die Inklusion ist nach
\FormulaRefAuto{InclusionOrderOnFamily}[thm] eine Teilordnung:
\[
  \PartOrd{\mathcal I_{A,\star},\subseteq}.
\]
Nach
\FormulaRefAuto{FinitePosetMinimalElement}[thm] besitzt sie bezüglich
\(\subseteq\) ein minimales Mitglied \(K\).  Ist \(I\) ein nichtleeres
zweiseitiges Ideal, so ist \(K\cap I\) nach
\FormulaRefAuto{SemigroupIdealIntersection}[thm] wieder ein nichtleeres
zweiseitiges Ideal und liegt in \(K\).  Die Minimalität von \(K\) liefert
\(K\cap I=K\), also \(K\subseteq I\).  Das minimale Ideal ist damit sogar
das kleinste Ideal.
\end{proof}

\FormulaThmDeltaK[Eindeutigkeit eines kleinsten nichtleeren Ideals]{%
  \begin{aligned}[t]
    &\SemigroupAlg{A}{\star}\dsep
      \mathsf{Id}(K;A,\star)\dsep
      \forall I\,\bigl(\mathsf{Id}(I;A,\star)\rightarrow K\subseteq I\bigr)
      \dsep{}\\[-2pt]
    &\mathsf{Id}(J;A,\star)\dsep
      \forall I\,\bigl(\mathsf{Id}(I;A,\star)\rightarrow J\subseteq I\bigr)
      \vdash K=J
  \end{aligned}%
}{FiniteSemigroupLeastIdealUniqueness}{%
  \DeltaRow{Mengen}{A\dsep I\dsep J\dsep K}
  \DeltaRow{Funktionen}{\star}
}
\begin{proof}
Die Universalbedingung für \(K\), angewandt auf das Ideal \(J\), ergibt
\(K\subseteq J\).  Umgekehrt liefert die Universalbedingung für \(J\),
angewandt auf \(K\), die Inklusion \(J\subseteq K\).  Also ist \(K=J\).
\end{proof}

\FormulaThmDeltaK[Minimalidealsatz für endliche Halbgruppen]{%
  \begin{aligned}[t]
    &\FiniteSemigroupAlg{A}{\star}\dsep A\neq\varnothing\vdash{}\\[-2pt]
    &\exists!K\,
      \Bigl(\mathsf{Id}(K;A,\star)\land
        \forall I\,\bigl(\mathsf{Id}(I;A,\star)\rightarrow K\subseteq I\bigr)
      \Bigr)
  \end{aligned}%
}{FiniteSemigroupLeastIdeal}{%
  \DeltaRow{Mengen}{A\dsep I\dsep J\dsep K}
  \DeltaRow{Funktionen}{\star}
}
\begin{proof}
Die Existenz folgt aus
\FormulaRefAuto{FiniteSemigroupLeastIdealExistence}[thm].  Sind \(K\) und
\(J\) zwei Mengen mit der angegebenen Eigenschaft, so liefert
\FormulaRefAuto{FiniteSemigroupLeastIdealUniqueness}[thm] die Gleichheit
\(K=J\).
\end{proof}

\FormulaDefDeltaK[Minimalideal einer nichtleeren endlichen Halbgruppe]{%
  \begin{aligned}[t]
    &\FiniteSemigroupAlg{A}{\star}\dsep A\neq\varnothing\vdash{}\\[-2pt]
    &K(A,\star)\coloneqq
      \iota K\,
      \Bigl(\mathsf{Id}(K;A,\star)\land
        \forall I\,\bigl(\mathsf{Id}(I;A,\star)\rightarrow K\subseteq I\bigr)
      \Bigr)
  \end{aligned}%
}{FiniteSemigroupLeastIdealDef}{%
  \DeltaRow{Mengen}{A\dsep I\dsep K}
  \DeltaRow{Funktionen}{\star}
  \DeltaRow{\textbf{Neues Symbol}}{}
  \DeltaPrem{Minimalideal}{K(A,\star)}
}

\FormulaThmDeltaK[Charakterisierung des Minimalideals]{%
  \begin{aligned}[t]
    &\FiniteSemigroupAlg{A}{\star}\dsep A\neq\varnothing\vdash{}\\[-2pt]
    &\mathsf{Id}\bigl(K(A,\star);A,\star\bigr)\land
      \forall I\,\bigl(
        \mathsf{Id}(I;A,\star)\rightarrow K(A,\star)\subseteq I
      \bigr)
  \end{aligned}%
}{FiniteSemigroupLeastIdealCharacterization}{%
  \DeltaRow{Mengen}{A\dsep I}
  \DeltaRow{Funktionen}{\star}
}
\begin{proof}
Der Existenz- und Eindeutigkeitssatz liefert das kleinste nichtleere Ideal.
Die beiden angezeigten Eigenschaften sind genau die definierenden
Eigenschaften der dadurch eingeführten Bezeichnung \(K(A,\star)\).
\end{proof}

%============================================================
\ProofWideReasonsAtRight
\chapter{Multiplikationstafeln und Potenzhalbgruppen}
%============================================================

\section{Identifikation durch Multiplikationstafeln}

\subsection{Nummerierte Multiplikationstafeln}

Ist \(\FiniteSemigroupAlg{A}{\star}\), so liefert die
Bijektionscharakterisierung der Endlichkeit eine natürliche Zahl \(n\) und
eine Nummerierung
\[
  e\colon\NatSeg{n}\bij A.
\]
Dabei beginnt die Nummerierung mit dem Index \(0\).  Die Verknüpfung auf
\(A\) lässt sich mit \(e\) auf den Peano-Anfangsabschnitt
\(\NatSeg{n}\) zurücktransportieren.

\FormulaDefDeltaK[Nummerierte Multiplikationstafel einer endlichen
Halbgruppe]{%
  \begin{aligned}[t]
    &\FiniteSemigroupAlg{A}{\star}\dsep n\in\mathbb N\dsep
      e\colon\NatSeg{n}\bij A\vdash{}\\[-2pt]
    &m_{\star,e}\colon
      \NatSeg{n}\times\NatSeg{n}\to\NatSeg{n},\\[-2pt]
    &m_{\star,e}(i,j)
      \coloneqq e^{-1}\!\bigl(e(i)\star e(j)\bigr)
      \qquad\bigl((i,j)\in\NatSeg{n}\times\NatSeg{n}\bigr)
  \end{aligned}%
}{FiniteSemigroupMultiplicationTableDef}{%
  \DeltaRow{Mengen}{A\dsep n\dsep i\dsep j}
  \DeltaRow{Funktionen}{e\dsep\star\dsep m_{\star,e}}
  \DeltaRow{\textbf{Axiome}}{}
  \DeltaPrem{Endliche Halbgruppe}{\FiniteSemigroupAlg{A}{\star}}
  \DeltaPrem{Länge der Nummerierung}{n\in\mathbb N}
  \DeltaPrem{Nummerierung}{e\colon\NatSeg{n}\bij A}
  \DeltaPrem{Indizes}{i,j\in\NatSeg{n}}
  \DeltaRow{\textbf{Neues Symbol}}{}
  \DeltaPrem{Multiplikationstafel}{m_{\star,e}}
}
\begin{tabproofwide}
  \proofstepwidestar[1]{\FiniteSemigroupAlg{A}{\star}}{\rA}
  \proofstepwidestar[2]{n\in\mathbb N}{\rA}
  \proofstepwidestar[3]{e\colon\NatSeg{n}\bij A}{\rA}
  \proofstepwidestar[4]{%
    (i,j)\in\NatSeg{n}\times\NatSeg{n}}{\rA}
  \proofstepwidestar[1]{\SemigroupAlg{A}{\star}}{%
    \FormulaRefAuto{FiniteSemigroupBaseAxiom}[ax]{1}}
  \proofstepwidestar[4]{i\in\NatSeg{n}}{%
    \FormulaRefAuto{(a,b)\in A\times B\vdash a\in A}{4}}
  \proofstepwidestar[4]{j\in\NatSeg{n}}{%
    \FormulaRefAuto{(a,b)\in A\times B\vdash b\in B}{4}}
  \proofstepwidestar[3,4]{e(i)\in A}{%
    \FormulaRefAuto{F\colon A\bij B\dsep x\in A\vdash F(x)\in B}%
      [thm]{3,6}}
  \proofstepwidestar[3,4]{e(j)\in A}{%
    \FormulaRefAuto{F\colon A\bij B\dsep x\in A\vdash F(x)\in B}%
      [thm]{3,7}}
  \proofstepwidestar[1,3,4]{e(i)\star e(j)\in A}{%
    \FormulaRefAuto{SemigroupClosureAxiom}[ax]{5,8,9}}
  \proofstepwidestar[1,2,3,4]{%
    e^{-1}\!\bigl(e(i)\star e(j)\bigr)\in\NatSeg{n}}{%
    \FormulaRefAuto{F\colon A\to B\dsep x\in A\vdash F(x)\in B}[thm]{%
      \FormulaRefAuto{F\colon A\bij B\vdash F^{-1}\colon B\to A}%
        [thm]{3},10}}
  \proofstepwidestar[1,2,3]{%
    \forall (i,j)\in\NatSeg{n}\times\NatSeg{n}\,
      e^{-1}\!\bigl(e(i)\star e(j)\bigr)\in\NatSeg{n}}{%
    \rUI{\rRI{4,11}}}
\end{tabproofwide}

Die Abgeschlossenheit von \(\star\) und die Bijektivität von \(e\)
stellen sicher, dass jeder Tafeleintrag wieder in \(\NatSeg{n}\) liegt.
Damit ist \(m_{\star,e}\) eine binäre Operation auf
\(\NatSeg{n}\).  Ihre Werte sind genau die Einträge der
Multiplikationstafel: Der erste Index bezeichnet die Zeile, der zweite die
Spalte.

\FormulaThmDeltaK[Nummerierte Tafeln transportieren die
Halbgruppenstruktur]{%
  \begin{aligned}[t]
    &\FiniteSemigroupAlg{A}{\star}\dsep n\in\mathbb N\dsep
      e\colon\NatSeg{n}\bij A\vdash{}\\[-2pt]
    &\FiniteSemigroupAlg{\NatSeg{n}}{m_{\star,e}}\land
      \FiniteSemigroupIso{e}{\NatSeg{n},m_{\star,e}}{A,\star}
  \end{aligned}%
}{FiniteSemigroupMultiplicationTableTransport}{%
  \DeltaRow{Mengen}{A\dsep n\dsep i\dsep j\dsep k}
  \DeltaRow{Funktionen}{e\dsep\star\dsep m_{\star,e}}
}
\begin{proof}
Für alle \(i,j\in\NatSeg n\) liefert die Definition zusammen mit der
Bijektivität von \(e\) die Transportgleichung
\[
  e\bigl(m_{\star,e}(i,j)\bigr)=e(i)\star e(j).
\]
Sind zusätzlich \(k\in\NatSeg n\), so haben die beiden Klammerungen unter
\(e\) die Bilder
\[
\begin{aligned}
e\bigl(m_{\star,e}(m_{\star,e}(i,j),k)\bigr)
  &=(e(i)\star e(j))\star e(k),\\
e\bigl(m_{\star,e}(i,m_{\star,e}(j,k))\bigr)
  &=e(i)\star(e(j)\star e(k)).
\end{aligned}
\]
Sie sind wegen der Assoziativität von \(\star\) gleich.  Die Injektivität
von \(e\) ergibt deshalb die Assoziativität von \(m_{\star,e}\).  Der
Anfangsabschnitt \(\NatSeg n\) ist endlich, und die Transportgleichung zeigt,
dass die gegebene Bijektion \(e\) ein Homomorphismus ist.  Damit folgen beide
Behauptungen.
\end{proof}

\subsection{Das exakte Umbenennungskriterium}

Auf einem bereits nummerierten Träger ist eine bijektive Abbildung genau
dann ein Isomorphismus, wenn sie alle Tafeleinträge erhält.

\FormulaThmDeltaK[Umbenennungskriterium für endliche
Halbgruppentafeln]{%
  \begin{aligned}[t]
    &n\in\mathbb N\dsep
      \FiniteSemigroupAlg{\NatSeg{n}}{\star}\dsep
      \FiniteSemigroupAlg{\NatSeg{n}}{\diamond}\dsep
      \pi\colon\NatSeg{n}\bij\NatSeg{n}\vdash{}\\[-2pt]
    &\FiniteSemigroupIso{\pi}
      {\NatSeg{n},\star}{\NatSeg{n},\diamond}
      \leftrightarrow
      \forall i,j\in\NatSeg{n}\,
        \bigl(\pi(i\star j)=\pi(i)\diamond\pi(j)\bigr)
  \end{aligned}%
}{FiniteSemigroupTableRelabelCriterion}{%
  \DeltaRow{Mengen}{n\dsep i\dsep j}
  \DeltaRow{Funktionen}{\pi\dsep\star\dsep\diamond}
}
\begin{proof}
Die Träger, ihre Halbgruppenstrukturen und die Bijektivität von \(\pi\) sind
bereits vorausgesetzt.  Nach den Definitionen des endlichen
Halbgruppenisomorphismus und des Halbgruppenhomomorphismus bleibt daher genau
die Operationserhaltung zu prüfen.  Diese lautet punktweise
\[
  \pi(i\star j)=\pi(i)\diamond\pi(j)
  \qquad(i,j\in\NatSeg n),
\]
und ist damit zugleich notwendig und hinreichend.
\end{proof}

\subsection{Umbenennungen und Isomorphieorbits}

Die Tafellogik lässt sich ohne Rekursion formulieren.  Für eine Menge
\(S\) bezeichnen wir mit \(\operatorname{Assoc}(S)\) die assoziativen
binären Operationen auf \(S\) und mit \(\operatorname{Sym}(S)\) ihre
Selbstbijektionen.

\FormulaDefDeltaK[Assoziative binäre Operationen auf einer Menge]{%
  \operatorname{Assoc}(S)\coloneqq
  \bigl\{m\in\powerset((S\times S)\times S)\bigm|
    m\colon S\times S\to S\land\SemigroupAlg{S}{m}\bigr\}%
}{FiniteSemigroupAssociativeOperationsDef}{%
  \DeltaRow{Mengen}{S}
  \DeltaRow{Funktionen}{m}
  \DeltaRow{\textbf{Neues Symbol}}{}
  \DeltaPrem{Assoziative Operationen}{\operatorname{Assoc}(S)}
}

\FormulaDefDeltaK[Permutationen einer Menge]{%
  \operatorname{Sym}(S)\coloneqq
  \bigl\{\pi\in\powerset(S\times S)\bigm|\pi\colon S\bij S\bigr\}%
}{FiniteSemigroupSymmetricSetDef}{%
  \DeltaRow{Mengen}{S}
  \DeltaRow{Funktionen}{\pi}
  \DeltaRow{\textbf{Neues Symbol}}{}
  \DeltaPrem{Permutationen}{\operatorname{Sym}(S)}
}


\FormulaThmDeltaK[Zugriff auf die Halbgruppenstruktur einer assoziativen
Operation]{%
  m\in\operatorname{Assoc}(S)\vdash\SemigroupAlg{S}{m}%
}{FiniteSemigroupAssociativeOperationSemigroup}{%
  \DeltaRow{Mengen}{S}
  \DeltaRow{Funktionen}{m}
}
\begin{tabproofwide}
  \proofstepwidestar[1]{m\in\operatorname{Assoc}(S)}{\rA}
  \proofstepwidestar[1]{\SemigroupAlg{S}{m}}{%
    \FormulaRefAuto{FiniteSemigroupAssociativeOperationsDef}[def]{1}}
\end{tabproofwide}

\FormulaThmDeltaK[Zugriff auf eine Selbstbijektion]{%
  \pi\in\operatorname{Sym}(S)\vdash\pi\colon S\bij S%
}{FiniteSemigroupSymmetricSetBijection}{%
  \DeltaRow{Mengen}{S}
  \DeltaRow{Funktionen}{\pi}
}
\begin{tabproofwide}
  \proofstepwidestar[1]{\pi\in\operatorname{Sym}(S)}{\rA}
  \proofstepwidestar[1]{\pi\colon S\bij S}{%
    \FormulaRefAuto{FiniteSemigroupSymmetricSetDef}[def]{1}}
\end{tabproofwide}

\FormulaDefDeltaK[Umbenannte Operation]{%
  \begin{aligned}[t]
    &m\in\operatorname{Assoc}(S)\dsep
      \pi\in\operatorname{Sym}(S)\vdash{}\\[-2pt]
    &m^\pi\colon S\times S\to S,\\[-2pt]
    &m^\pi(u,v)\coloneqq
      \pi\!\left(m\bigl(\pi^{-1}(u),\pi^{-1}(v)\bigr)\right)
      \qquad(u,v\in S)
  \end{aligned}%
}{FiniteSemigroupRelabelledOperationDef}{%
  \DeltaRow{Mengen}{S\dsep u\dsep v}
  \DeltaRow{Funktionen}{m\dsep\pi\dsep m^\pi}
  \DeltaRow{\textbf{Neues Symbol}}{}
  \DeltaPrem{Umbenannte Operation}{m^\pi}
}

\FormulaThmDeltaK[Die Umbenennung ist eine Wirkung auf der Menge der
assoziativen Operationen]{%
  \begin{aligned}[t]
    &m\in\operatorname{Assoc}(S)\dsep
      \pi,\sigma\in\operatorname{Sym}(S)\vdash{}\\[-2pt]
    &\Id_S\in\operatorname{Sym}(S)\land
      \sigma\circ\pi\in\operatorname{Sym}(S)\land{}\\[-2pt]
    &m^\pi\in\operatorname{Assoc}(S)\land
      m^{\Id_S}=m\land
      (m^\pi)^\sigma=m^{\sigma\circ\pi}
  \end{aligned}%
}{FiniteSemigroupRelabellingAction}{%
  \DeltaRow{Mengen}{S\dsep u\dsep v\dsep w}
  \DeltaRow{Funktionen}{m\dsep\pi\dsep\sigma}
}
\begin{proof}
Da \(m\colon S\times S\to S\) gilt und \(\pi\) eine Selbstbijektion ist,
definiert die angegebene Formel zunächst wieder eine binäre Operation auf
\(S\).  Außerdem gelten
\[
  \Id_S\in\operatorname{Sym}(S),
  \qquad
  \sigma\circ\pi\in\operatorname{Sym}(S),
\]
weil die Identität und die Komposition von Bijektionen bijektiv sind.
Für \(u,v\in S\) liefert die Definition zunächst
\[
  \pi^{-1}\!\bigl(m^\pi(u,v)\bigr)
  =m\bigl(\pi^{-1}(u),\pi^{-1}(v)\bigr).
\]
Damit folgt für \(u,v,w\in S\)
\[
\begin{aligned}
m^\pi(m^\pi(u,v),w)
 &=\pi\!\left(
   m\bigl(m(\pi^{-1}(u),\pi^{-1}(v)),\pi^{-1}(w)\bigr)
   \right)\\
 &=\pi\!\left(
   m\bigl(\pi^{-1}(u),m(\pi^{-1}(v),\pi^{-1}(w))\bigr)
   \right)\\
 &=m^\pi(u,m^\pi(v,w)).
\end{aligned}
\]
Also ist \(m^\pi\in\operatorname{Assoc}(S)\); damit ist nun auch
\((m^\pi)^\sigma\) definiert.  Für alle \(u,v\in S\) gilt punktweise
\[
  m^{\Id_S}(u,v)=m(u,v)
\]
und wegen
\((\sigma\circ\pi)^{-1}=\pi^{-1}\circ\sigma^{-1}\)
\[
  (m^\pi)^\sigma(u,v)
  =\sigma\!\left(\pi\!\left(
      m\bigl(\pi^{-1}(\sigma^{-1}(u)),
             \pi^{-1}(\sigma^{-1}(v))\bigr)
    \right)\right)
  =m^{\sigma\circ\pi}(u,v).
\]
Die Funktionsextensionalität liefert daraus
\(m^{\Id_S}=m\) und
\((m^\pi)^\sigma=m^{\sigma\circ\pi}\).
\end{proof}

\FormulaDefDeltaK[Graph der Umbenennungswirkung]{%
  \begin{aligned}[t]
    &m\in\operatorname{Assoc}(S)\vdash{}\\[-2pt]
    &\operatorname{Relab}_S(\pi;m,d)\coloneqq
      \pi\in\operatorname{Sym}(S)\land
      d\in\operatorname{Assoc}(S)\land{}\\[-2pt]
    &\hspace{4em}
      \forall x,y\in S\,
        \bigl(\pi(m(x,y))=d(\pi(x),\pi(y))\bigr)
  \end{aligned}%
}{FiniteSemigroupRelabelGraphDef}{%
  \DeltaRow{Mengen}{S\dsep x\dsep y}
  \DeltaRow{Funktionen}{m\dsep d\dsep\pi}
  \DeltaRow{\textbf{Neues Symbol}}{}
  \DeltaPrem{Umbenennungsrelation}{\operatorname{Relab}_S(\pi;m,d)}
}

\FormulaThmDeltaK[Die Umbenennungsrelation ist der Graph der Wirkung]{%
  \begin{aligned}[t]
    &m,d\in\operatorname{Assoc}(S)\dsep
      \pi\in\operatorname{Sym}(S)\vdash{}\\[-2pt]
    &\operatorname{Relab}_S(\pi;m,d)\leftrightarrow d=m^\pi
  \end{aligned}%
}{FiniteSemigroupRelabelGraphCharacterization}{%
  \DeltaRow{Mengen}{S\dsep u\dsep v\dsep x\dsep y}
  \DeltaRow{Funktionen}{m\dsep d\dsep\pi}
}
\begin{proof}
Aus der Homomorphiegleichung der Relation folgt nach Einsetzen von
\(x=\pi^{-1}(u)\) und \(y=\pi^{-1}(v)\) punktweise
\[
  d(u,v)=\pi\!\left(m(\pi^{-1}(u),\pi^{-1}(v))\right)=m^\pi(u,v).
\]
Umgekehrt liefert \(d=m^\pi\) nach Einsetzen von \(u=\pi(x)\) und
\(v=\pi(y)\) unmittelbar
\(d(\pi(x),\pi(y))=\pi(m(x,y))\).
\end{proof}

\FormulaDefDeltaK[Umbenennungsorbit einer assoziativen Operation]{%
  \begin{aligned}[t]
    &m\in\operatorname{Assoc}(S)\vdash{}\\[-2pt]
    &\operatorname{Orb}_S(m)\coloneqq
      \bigl\{d\in\operatorname{Assoc}(S)\bigm|
        \exists\pi\in\operatorname{Sym}(S)\,
          \operatorname{Relab}_S(\pi;m,d)\bigr\}
  \end{aligned}%
}{FiniteSemigroupRelabelOrbitDef}{%
  \DeltaRow{Mengen}{S}
  \DeltaRow{Funktionen}{m\dsep d\dsep\pi}
  \DeltaRow{\textbf{Neues Symbol}}{}
  \DeltaPrem{Umbenennungsorbit}{\operatorname{Orb}_S(m)}
}

\FormulaDefDeltaK[Automorphismengruppe einer assoziativen Operation]{%
  \begin{aligned}[t]
    &m\in\operatorname{Assoc}(S)\vdash{}\\[-2pt]
    &\operatorname{Aut}(S,m)\coloneqq
      \bigl\{\pi\in\operatorname{Sym}(S)\bigm|m^\pi=m\bigr\}
  \end{aligned}%
}{FiniteSemigroupAutomorphismGroupDef}{%
  \DeltaRow{Mengen}{S}
  \DeltaRow{Funktionen}{m\dsep\pi}
  \DeltaRow{\textbf{Neues Symbol}}{}
  \DeltaPrem{Automorphismengruppe}{\operatorname{Aut}(S,m)}
}

Diese Menge ist tatsächlich eine Untergruppe von
\(\operatorname{Sym}(S)\).  Aus dem Wirkungsgesetz folgen
\[
  m^{\Id_S}=m,\qquad
  (m^\pi)^\sigma=m^{\sigma\circ\pi}.
\]
Damit enthält \(\operatorname{Aut}(S,m)\) die Identität und ist unter
Komposition abgeschlossen.  Ist \(m^\pi=m\), so liefert die Umbenennung mit
\(\pi^{-1}\)
\[
  m^{\pi^{-1}}
  =(m^\pi)^{\pi^{-1}}
  =m^{\pi^{-1}\circ\pi}
  =m;
\]
also enthält die Menge mit jedem Element auch dessen Inverses.

\FormulaThmDeltaK[Umbenennungsorbits sind die Isomorphieklassen auf dem
festen Träger]{%
  \begin{aligned}[t]
    &\Finite{S}\dsep m,d\in\operatorname{Assoc}(S)\vdash{}\\[-2pt]
    &d\in\operatorname{Orb}_S(m)
      \leftrightarrow
      \exists f\,\FiniteSemigroupIso{f}{S,m}{S,d}
  \end{aligned}%
}{FiniteSemigroupRelabelOrbitsAreIsomorphismClasses}{%
  \DeltaRow{Mengen}{S\dsep x\dsep y}
  \DeltaRow{Funktionen}{m\dsep d\dsep f\dsep\pi}
}
\begin{proof}
Ist \(d\in\operatorname{Orb}_S(m)\), so gibt es eine Permutation \(\pi\)
mit \(d=m^\pi\).  Die definierende Gleichung
\[
  \pi(m(x,y))=m^\pi(\pi(x),\pi(y))
\]
zeigt, dass \(\pi\) ein Isomorphismus von \((S,m)\) nach \((S,d)\) ist;
wegen der vorausgesetzten Endlichkeit ist dies auch ein Isomorphismus
endlicher Halbgruppen.  Umgekehrt ist jeder solche Isomorphismus \(f\) eine
Permutation von \(S\), und seine Homomorphiegleichung ergibt nach der
vorherigen Charakterisierung \(d=m^f\).  Also liegt \(d\) im Orbit von
\(m\).
\end{proof}

Nun seien \(e_A\colon\NatSeg{n}\bij A\) und
\(e_B\colon\NatSeg{n}\bij B\) Nummerierungen zweier endlicher
Halbgruppen gleicher Trägergröße.  Auf ihre nummerierten Tafeln
angewandt ergibt der Satz das vollständige Identifikationskriterium
\[
\begin{aligned}
&\exists f\,
  \FiniteSemigroupIso{f}{A,\star}{B,\diamond}
\quad\Longleftrightarrow{}\\[-2pt]
&\exists\pi\Bigl(
  \pi\colon\NatSeg{n}\bij\NatSeg{n}\ \land
  \forall i,j\in\NatSeg{n}\,
  m_{\diamond,e_B}\bigl(\pi(i),\pi(j)\bigr)
    =\pi\bigl(m_{\star,e_A}(i,j)\bigr)
  \Bigr).
\end{aligned}
\]
Es muss also dieselbe Umbenennung auf die Zeilen, die Spalten und die
Tafeleinträge angewandt werden.  In der Hinrichtung ist
\(\pi=e_B^{-1}\circ f\circ e_A\); in der Rückrichtung gewinnt man den
Isomorphismus als \(f=e_B\circ\pi\circ e_A^{-1}\).  Dies beweist zugleich,
dass das Kriterium notwendig und hinreichend ist.

\subsection{Ein kanonischer Tafelcode}

Das Umbenennungskriterium liefert einen exakten, endlichen
Identifikationsalgorithmus.  Zunächst kodieren wir eine nummerierte Tafel
ohne mehrdeutige Ziffernverkettung.

\FormulaDefDeltaK[Zeilenweise gelesenes Tafeltupel]{%
  \begin{aligned}[t]
    &n\in\mathbb N\dsep
      m\colon\NatSeg n\times\NatSeg n\to\NatSeg n\vdash{}\\[-2pt]
    &w_n(m)\coloneqq
      \Bigl(\bigl(m(i,j)\bigr)_{j\in\NatSeg n}\Bigr)_{i\in\NatSeg n}
  \end{aligned}%
}{FiniteSemigroupTableTupleDef}{%
  \DeltaRow{Mengen}{n\dsep i\dsep j}
  \DeltaRow{Funktionen}{m\dsep w_n}
  \DeltaRow{\textbf{Neues Symbol}}{}
  \DeltaPrem{Tafeltupel}{w_n(m)}
}

\FormulaThmDeltaK[Das Tafeltupel bestimmt die Tafel]{%
  \begin{aligned}[t]
    &n\in\mathbb N\dsep
      m,d\colon\NatSeg n\times\NatSeg n\to\NatSeg n\dsep
      w_n(m)=w_n(d)\vdash m=d
  \end{aligned}%
}{FiniteSemigroupTableTupleInjective}{%
  \DeltaRow{Mengen}{n\dsep i\dsep j}
  \DeltaRow{Funktionen}{m\dsep d\dsep w_n}
}
\begin{proof}
Die Gleichheit der äußeren Tupel liefert für jedes \(i\in\NatSeg n\) die
Gleichheit der zugehörigen Zeilentupel.  Deren \(j\)-te Komponenten stimmen
für jedes \(j\in\NatSeg n\) überein, also \(m(i,j)=d(i,j)\).  Nach der
Funktionsextensionalität folgt \(m=d\).
\end{proof}

Die äußere Folge durchläuft die Zeilen, die innere jeweils die Spalten.  Wir
identifizieren dieses verschachtelte Tupel mit seinem flachen Tupel der Länge
\(n^2\).  Es ist keine Dezimalzahl; dadurch bleibt die Schreibweise auch für
\(n\geq11\) eindeutig.  Beispielsweise wird die
Tafel
\[
\begin{array}{c|cc}
\star&0&1\\\hline
0&0&0\\
1&1&1
\end{array}
\]
ohne Zeilen- und Spaltenköpfe als \(w_2(m)=(0,0,1,1)\) gelesen.  Da alle
Nummerierungen in die Kandidatenmenge eingehen, ist der folgende Code von der
Wahl einer Ausgangsnummerierung unabhängig.

\FormulaDefDeltaK[Kanonischer Tafelcode einer endlichen Halbgruppe]{%
  \begin{aligned}[t]
    &\FiniteSemigroupAlg{A}{\star}\vdash{}\\[-2pt]
    &\operatorname{can}(A,\star)\coloneqq
      \left(
        \tCard(A),
        \left\{
          w_{\tCard(A)}\!\left(m_{\star,e}\right)
          \;\middle|\;
          e\colon\NatSeg{\tCard(A)}\bij A
        \right\}
      \right)
  \end{aligned}%
}{FiniteSemigroupCanonicalTableCodeDef}{%
  \DeltaRow{Mengen}{A}
  \DeltaRow{Funktionen}{\star\dsep e\dsep m_{\star,e}\dsep w_{\tCard(A)}}
  \DeltaRow{\textbf{Neues Symbol}}{}
  \DeltaPrem{Kanonischer Tafelcode}{\operatorname{can}(A,\star)}
}

Die Menge der Nummerierungen ist nichtleer und endlich.  Folglich ist auch
die im Code gespeicherte Menge der Tafeltupel nichtleer und endlich.  Die
erste Komponente hält die Trägergröße ausdrücklich fest.

\FormulaThmDeltaK[Der kanonische Tafelcode ist vollständig]{%
  \begin{aligned}[t]
    &\FiniteSemigroupAlg{A}{\star}\dsep
      \FiniteSemigroupAlg{B}{\diamond}\vdash{}\\[-2pt]
    &\bigl(\exists f\,\FiniteSemigroupIso{f}{A,\star}{B,\diamond}\bigr)
      \leftrightarrow
      \operatorname{can}(A,\star)=\operatorname{can}(B,\diamond)
  \end{aligned}%
}{FiniteSemigroupCanonicalTableCodeComplete}{%
  \DeltaRow{Mengen}{A\dsep B}
  \DeltaRow{Funktionen}{\star\dsep\diamond\dsep f}
}
\begin{proof}
Sei zunächst \(f\) ein Isomorphismus von \((A,\star)\) nach
\((B,\diamond)\).  Dann ist \(\tCard(A)=\tCard(B)=n\), und die Zuordnung
\(e\mapsto f\circ e\) ist eine Bijektion zwischen ihren Nummerierungen.  Die
Homomorphie von \(f\) zeigt für jede Nummerierung \(e\) punktweise
\[
  m_{\diamond,f\circ e}(i,j)=m_{\star,e}(i,j).
\]
Beide Kandidatenmengen und damit beide Codes stimmen also überein.

Umgekehrt erzwingt die Gleichheit der Codes zunächst die gleiche Trägergröße
\(n\).  Wähle irgendeine Nummerierung \(e_A\) von \(A\).  Ihr Tafeltupel
gehört wegen der Gleichheit der zweiten Codekomponenten auch zum Code von
\(B\); also gibt es eine Nummerierung \(e_B\) mit
\[
  w_n(m_{\star,e_A})=w_n(m_{\diamond,e_B}).
\]
Nach
\FormulaRefAuto{FiniteSemigroupTableTupleInjective}[thm] gilt
\(m_{\star,e_A}=m_{\diamond,e_B}\).  Der Strukturtransport zeigt nun, dass
\(e_B\circ e_A^{-1}\) ein Isomorphismus von \((A,\star)\) nach
\((B,\diamond)\) ist.
\end{proof}

Für die Berechnung genügt weiterhin eine feste Nummerierung
\(e\colon\NatSeg n\bij A\).  Jede andere Nummerierung hat die Form
\(e\circ\rho\) mit \(\rho\in\operatorname{Sym}(\NatSeg n)\), und
\[
  m_{\star,e\circ\rho}=m_{\star,e}^{\rho^{-1}}.
\]
Die Kandidatenmenge des Codes ist daher genau die Menge der Tafeltupel im
Umbenennungsorbit einer beliebigen Ausgangstafel.  Rechnerisch darf man diese
endliche Menge auf ihr kleinstes Element in der üblichen zeilenweisen
lexikographischen Ordnung komprimieren: Zwei Orbits haben genau dann dasselbe
Minimum, wenn sie derselbe Orbit sind.  Diese Ordnung dient nur der Wahl eines
kurzen Vertreters; der formal definierte Code selbst benötigt sie nicht.

Für einen Träger mit \(n\) Elementen gibt es genau \(n!\) Permutationen und
höchstens \(n!\) verschiedene umbenannte Tafeln; je Permutation sind \(n^2\)
Tafeleinträge zu prüfen.  Bei
zwei beziehungsweise drei Elementen sind dies nur \(2\) beziehungsweise
\(6\) Umbenennungen.  Als schnelle Vorfilter kann man etwa die Anzahl der
idempotenten Elemente, Kommutativität, die Existenz eines Nullelements,
Kürzbarkeit und die in Band~22 eingeführten Teilbarkeitsrelationen
vergleichen.  Verschiedene Werte schließen einen Isomorphismus aus;
übereinstimmende Werte ersetzen den kanonischen Tafelcode im Allgemeinen
nicht.

\begin{remark}[Grenze des Tafelverfahrens]
Das Tafelverfahren identifiziert endliche Halbgruppen, sobald ihre
Multiplikationstafeln vorliegen.  Aus einem abstrakten Isomorphismus ihrer
Potenzhalbgruppen lässt sich jedoch nicht ohne Weiteres eine Umbenennung der
Grundträger ablesen: Einermengen sind rein multiplikativ nicht im
Allgemeinen ausgezeichnet.  Das abschließende Rekonstruktionskapitel macht
dies an der zweielementigen Nullhalbgruppe konkret: Dort kann sogar ein
Isomorphismus der Potenzhalbgruppe auf sich selbst eine Einermenge mit einer
zweielementigen Teilmenge vertauschen.  Der kanonische Tafelcode ist daher
ein vollständiger Isomorphietest für gegebene Tafeln, aber noch kein
allgemeiner Beweis der später formulierten Potenzhalbgruppenvermutung.
\end{remark}

%============================================================
\section{Nummerierung der Potenzhalbgruppe}
%============================================================

\subsection{Mächtigkeitsklassen und ihre Indexblöcke}

Für diesen Abschnitt seien \(N\in\mathbb N\), \(N\neq0\),
\(S_N=\NatSeg N\) und \(m\in\operatorname{Assoc}(S_N)\) fest.  Wir
schreiben \(x\star y=m(x,y)\).  Dann ist
\[
  \PowNE{S_N}
  =
  \{X\subseteq S_N\mid X\neq\varnothing\}.
\]
Die nichtleeren Teilmengen werden nach ihrer Mächtigkeit zerlegt:
\[
  \boxed{
  P_r
  \coloneqq
  \{X\in\PowNE{S_N}\mid |X|=r\}}
  \qquad(1\leq r\leq N).
\]
Damit gilt
\[
  \PowNE{S_N}
  =
  P_1\mathbin{\dot\cup}\cdots\mathbin{\dot\cup}P_N,
  \qquad
  |P_r|=\binom{N}{r}.
\]
Insbesondere ist
\[
  P_1=\bigl\{\{0\},\ldots,\{N-1\}\bigr\}.
\]

Für die nummerierte Darstellung wählen wir eine Bijektion
\[
  \nu\colon
  \NatSeg{2^N-1}
  \bij
  \PowNE{S_N}
\]
mit
\[
  \nu(i)=\{i\}
  \qquad(0\leq i<N).
\]
Dabei werden zuerst die Klassen \(P_1,\ldots,P_N\) nach wachsender
Mächtigkeit durchlaufen; innerhalb jeder Klasse ordnen wir die Teilmengen
lexikographisch nach ihren aufsteigend geschriebenen Elementtupeln.  Setze
\[
  b_r\coloneqq
  \sum_{q=1}^{r-1}\binom{N}{q}
  \qquad(1\leq r\leq N),
\]
wobei die leere Summe \(b_1=0\) ist.  Die zu \(P_r\) gehörenden
\emph{Indexblöcke} sind
\[
  \boxed{
  I_r
  \coloneqq
  \nu^{-1}[P_r]
  =
  \left\{
    b_r,\ldots,b_r+\binom{N}{r}-1
  \right\}.}
\]
Diese Notation trennt die beiden Ebenen sauber: \(P_r\) besteht aus
Teilmengen von \(S_N\), während \(I_r\) aus deren Nummern besteht.

Die Klassen \(P_r\) sind im Allgemeinen keine Unterhalbgruppen.  Für
\(X\in P_r\) und \(Y\in P_s\) gilt zwar
\[
  1\leq|X\star_{\mathcal P}Y|
  \leq\min\{N,rs\},
\]
aber verschiedene Paare \((x,y)\) können denselben Produktwert liefern.  Die
Mächtigkeit des Ergebnisses muss daher nicht allein von \(r\) und \(s\)
abhängen.

\begin{example}[Ein einziger Block mit drei verschiedenen Zielklassen]
Auf \(S_3=\{0,1,2\}\) betrachte die assoziative Tafel
\[
\begin{array}{c|ccc}
\star&0&1&2\\\hline
0&0&0&0\\
1&0&1&2\\
2&0&2&0
\end{array}.
\]
Hier ist \(1\) neutral, \(0\) absorbierend und \(2\star2=0\).  Die
Assoziativität ist schnell geprüft: Tripel mit \(0\) kollabieren zu \(0\),
\(1\) kann aus jeder Klammerung entfernt werden, und im einzigen verbleibenden
Fall \((2,2,2)\) ergeben beide Klammerungen \(0\).  Für die drei
zweielementigen Teilmengen ergeben sich
\[
\begin{aligned}
  \{0,2\}\star_{\mathcal P}\{0,2\}&=\{0\}\in P_1,\\
  \{0,1\}\star_{\mathcal P}\{0,1\}&=\{0,1\}\in P_2,\\
  \{1,2\}\star_{\mathcal P}\{1,2\}&=\{0,1,2\}\in P_3.
\end{aligned}
\]
Alle drei Produkte liegen im selben Eingabeblock \(P_2\times P_2\), landen
aber in drei verschiedenen Mächtigkeitsklassen.  Eine Blockposition legt
also im Allgemeinen nur fest, welche Teilrechteckgröße gelesen wird, nicht
die Größe des Ergebnisses.
\end{example}


%============================================================
\subsection{Teilmengenprodukte als Tafelblöcke}
%============================================================

Für nichtleere Teilmengen \(X,Y\subseteq S_N\) lautet das Produkt
\[
  X\star_{\mathcal P}Y
  =
  \{m(x,y)\mid x\in X,\ y\in Y\}.
\]
Über \(\nu\) wird daraus die nummerierte Tafel
\[
  \boxed{
  \begin{aligned}
  m_{\mathcal P}\colon{}
    &\NatSeg{2^N-1}\times\NatSeg{2^N-1}
      \longrightarrow\NatSeg{2^N-1},\\[-2pt]
  m_{\mathcal P}(p,q)
    &\coloneqq
      \nu^{-1}\!\left(
        \{m(x,y)\mid x\in\nu(p),\ y\in\nu(q)\}
      \right).
  \end{aligned}}
\]
Die Wertemenge ist nichtleer, weil \(\nu(p)\) und \(\nu(q)\) nichtleer
sind.  Somit ist \(\nu^{-1}\) anwendbar.  Genauer ist
\(m_{\mathcal P}=m_{\star_{\mathcal P},\nu}\); die Assoziativität folgt
also aus dem Satz über endliche Potenzhalbgruppen und dem Strukturtransport
nummerierter Tafeln.
Sind \(p\in I_r\) und \(q\in I_s\), so liest man in der ursprünglichen
Tafel genau das \(r\times s\)-Teilrechteck mit den Zeilen aus \(\nu(p)\) und
den Spalten aus \(\nu(q)\).  Man entfernt doppelte Werte und nummeriert die
verbleibende Wertemenge mit \(\nu^{-1}\).

Zum Beispiel liefern
\[
  \nu(p)=\{i,j\},
  \qquad
  \nu(q)=\{k,\ell\}
\]
den Eintrag
\[
  m_{\mathcal P}(p,q)
  =
  \nu^{-1}\!\bigl(
    \{m(i,k),m(i,\ell),m(j,k),m(j,\ell)\}
  \bigr).
\]
Die ganze Potenzhalbgruppentafel zerfällt auf diese Weise sichtbar in die
Blöcke \(I_r\times I_s\).  Ihr Block \(I_1\times I_1\) ist exakt die
ursprüngliche Tafel, denn
\[
\begin{aligned}
  \{i\}\star_{\mathcal P}\{j\}&=\{m(i,j)\},\\
  m_{\mathcal P}(i,j)&=m(i,j)
    \qquad(i,j\in S_N=I_1).
\end{aligned}
\]


%============================================================
\section{Potenzhalbgruppe des Leitbeispiels}
%============================================================

Als Leitbeispiel des Haupttextes fixieren wir jetzt die endliche Halbgruppe
auf \(S_3=\{0,1,2\}\) mit der Multiplikationstafel
\[
\begin{array}{c|ccc}
\star&0&1&2\\\hline
0&2&0&1\\
1&0&1&2\\
2&1&2&0
\end{array}.
\]
Aus der Tafel liest man
\[
  x\star y\equiv x+y+2\pmod 3.
\]
Beide Klammerungen eines Dreifachprodukts sind daher zu
\(x+y+z+1\) modulo \(3\) kongruent.  Da ihre Werte in
\(S_3=\{0,1,2\}\) liegen, sind sie gleich; damit ist die Assoziativität ohne
Fallprüfung bewiesen.  Eine strukturelle Herleitung derselben Tafel aus der
Potenzfolge von \(0\) steht im rechnerischen Anhang.

Für \(S_3=\{0,1,2\}\) verwenden wir die Nummerierung
\[
\begin{array}{c|c|c}
p&\nu(p)&\text{Indexblock}\\\hline
0&\{0\}&I_1\\
1&\{1\}&I_1\\
2&\{2\}&I_1\\\hline
3&\{0,1\}&I_2\\
4&\{0,2\}&I_2\\
5&\{1,2\}&I_2\\\hline
6&\{0,1,2\}&I_3
\end{array}
\]
und damit
\[
  I_1=\{0,1,2\},
  \qquad
  I_2=\{3,4,5\},
  \qquad
  I_3=\{6\}.
\]

Zunächst ein Produkt aus dem Block \(I_1\times I_2\).  Für \(p=0\) und
\(q=4\) werden die Zeile \(0\) und die Spalten \(0,2\) der Grundtafel
gelesen:
\[
\begin{array}{c|cc}
&0&2\\\hline
0&2&1
\end{array}
\quad\Longrightarrow\quad
\{2,1\}=\{1,2\}=\nu(5).
\]
Also ist
\[
  m_{\mathcal P}(0,4)=5.
\]

Für das Produkt aus \(I_2\times I_2\) mit \(p=3\) und \(q=4\) wird ein
\(2\times2\)-Rechteck benötigt:
\[
\begin{array}{c|cc}
&0&2\\\hline
0&2&1\\
1&0&2
\end{array}
\quad\Longrightarrow\quad
\{2,1,0,2\}=\{0,1,2\}=\nu(6).
\]
Somit gilt
\[
  m_{\mathcal P}(3,4)=6.
\]
Als weiteres Beispiel ist
\[
  \{0,1\}\star_{\mathcal P}\{1\}=\{0,1\},
  \qquad\text{also}\qquad
  m_{\mathcal P}(3,1)=3.
\]

In diesem Beispiel hängt die Mächtigkeitsklasse des Ergebnisses tatsächlich
nur vom Zeilen- und Spaltenblock ab.  Die folgende kleine Blocktafel fasst
dieses Verhalten zusammen:
\[
\begin{array}{c|ccc}
&I_1&I_2&I_3\\\hline
I_1&I_1&I_2&I_3\\
I_2&I_2&I_3&I_3\\
I_3&I_3&I_3&I_3
\end{array}.
\]
Dies ist eine Besonderheit des Leitbeispiels und keine allgemeine Regel für
Mächtigkeitsklassen.

In der vollständigen nummerierten Tafel sind dieselben Blöcke nun durch
senkrechte und waagerechte Linien eingezeichnet:
\[
\boxed{
\begin{array}{cc|ccc|ccc|c}
\multicolumn{2}{c|}{}
  &\multicolumn{3}{c|}{I_1}
  &\multicolumn{3}{c|}{I_2}
  &\multicolumn{1}{c}{I_3}\\
\multicolumn{2}{c|}{m_{\mathcal P}}
  &0&1&2&3&4&5&6\\\hline
I_1&0&2&0&1&4&5&3&6\\
   &1&0&1&2&3&4&5&6\\
   &2&1&2&0&5&3&4&6\\\hline
I_2&3&4&3&5&6&6&6&6\\
   &4&5&4&3&6&6&6&6\\
   &5&3&5&4&6&6&6&6\\\hline
I_3&6&6&6&6&6&6&6&6
\end{array}}
\]
Der obere linke \(I_1\times I_1\)-Block ist die ursprüngliche
\(3\times3\)-Tafel.  Das durch \(6\) nummerierte Element
\(\nu(6)=S_3\) ist in dieser Potenzhalbgruppe absorbierend.  Für jede
nichtleere Teilmenge \(X\subseteq S_3\) gilt nämlich
\[
  X\star_{\mathcal P}S_3
  =
  S_3
  =
  S_3\star_{\mathcal P}X.
\]
Das lässt sich ohne Gruppentheorie direkt aus der Grundtafel lesen: Jede
vollständige Zeile enthält die Werte \(0,1,2\).  Wählt man ein
\(x\in X\), so ist bereits
\(\{x\}\star_{\mathcal P}S_3=S_3\).  Ebenso enthält jede vollständige
Spalte alle drei Werte, also
\(S_3\star_{\mathcal P}\{x\}=S_3\).
Entsprechend ist \(6\) für die nummerierte Operation \(m_{\mathcal P}\)
absorbierend.


%============================================================
\subsection{Kontrast: Potenzhalbgruppe der Nullhalbgruppe}
%============================================================

Betrachte auf \(S_2=\{0,1\}\) die Nullhalbgruppe \(m(a,b)=0\).  Mit
\[
  \nu(0)=\{0\},
  \qquad
  \nu(1)=\{1\},
  \qquad
  \nu(2)=\{0,1\}
\]
sind \(I_1=\{0,1\}\) und \(I_2=\{2\}\).  Für beliebige nichtleere
Teilmengen \(X,Y\subseteq S_2\) gilt
\[
  X\star_{\mathcal P}Y=\{0\}=\nu(0).
\]
Daher lautet die blockweise markierte Potenzhalbgruppentafel
\[
\begin{array}{cc|cc|c}
\multicolumn{2}{c|}{}
  &\multicolumn{2}{c|}{I_1}
  &\multicolumn{1}{c}{I_2}\\
\multicolumn{2}{c|}{m_{\mathcal P}}&0&1&2\\\hline
I_1&0&0&0&0\\
   &1&0&0&0\\\hline
I_2&2&0&0&0
\end{array}.
\]
Hier fällt sogar der Block \(I_2\times I_2\) in \(I_1\).  Das macht konkret
sichtbar, warum die Mächtigkeitsklassen zwar eine nützliche Gliederung der
Tafel, aber im Allgemeinen keine algebraisch abgeschlossenen Teile sind.

\subsection{Linksnull, Rechtsnull und ein Halbverband}

Zwei besonders kurze Beispiele zeigen, dass die Richtung der Grundoperation
im Mengenprodukt erhalten bleibt.  In der \emph{Linksnullhalbgruppe} gilt
\(x\star y=x\).  Daher ist für alle nichtleeren Teilmengen
\[
  X\star_{\mathcal P}Y=X.
\]
Die Potenzhalbgruppe ist also selbst eine Linksnullhalbgruppe, nun auf
\(2^N-1\) Elementen.  In der \emph{Rechtsnullhalbgruppe} mit
\(x\diamond y=y\) gilt entsprechend
\[
  X\diamond_{\mathcal P}Y=Y.
\]
Schon an einem einzigen Produkt ist damit sichtbar, welcher Faktor den Wert
bestimmt.

Als kommutativer Kontrast diene \(S_2=\{0,1\}\) mit
\(x\vee y=\max\{x,y\}\).  Schreibe
\[
  A=\{0\},\qquad B=\{1\},\qquad C=\{0,1\}.
\]
Dann lautet die Potenzhalbgruppentafel
\[
\begin{array}{c|ccc}
\vee_{\mathcal P}&A&B&C\\\hline
A&A&B&C\\
B&B&B&B\\
C&C&B&C
\end{array}.
\]
Alle drei Teilmengen sind idempotent, \(A\) ist neutral und \(B\)
absorbierend.  Insbesondere ist Idempotenz kein Merkmal, an dem sich
Einermengen in einer Potenzhalbgruppe allgemein erkennen lassen.


%============================================================
\chapter{Von der Potenzhalbgruppe zur Ausgangshalbgruppe}
%============================================================

Wo die Operationen aus dem Zusammenhang eindeutig sind, schreiben wir kurz
\(xy\) für das Grundprodukt und \(XY\) für das mengenweise Produkt
\(X\star_{\mathcal P}Y\); \(XY\) bezeichnet kein kartesisches Produkt.

\section{Einermengenkopie und Rekonstruktion}

Für jede Halbgruppe \((A,\star)\) ist die in
\FormulaRefAuto{SemigroupSingletonCopyDef}[def] definierte
Einermengenschicht
\[
  \operatorname{Sing}(A)=\bigl\{\{a\}\bigm|a\in A\bigr\}
\]
nach \FormulaRefAuto{SemigroupSingletonCopyIsomorphism}[thm] eine zu
\((A,\star)\) isomorphe Unterhalbgruppe ihrer Potenzhalbgruppe.

\FormulaThmDeltaK[Rekonstruktion bei erhaltener Einermengenschicht]{%
  \begin{aligned}[t]
    &\SemigroupAlg{A}{\star}\dsep
      \SemigroupAlg{B}{\diamond}\dsep{}\\[-2pt]
    &\SemigroupIso{F}
      {\PowNE{A},\star_{\mathcal P}}
      {\PowNE{B},\diamond_{\mathcal P}}\dsep
      F[\operatorname{Sing}(A)]=\operatorname{Sing}(B)\vdash{}\\[-2pt]
    &\exists f\,\SemigroupIso{f}{A,\star}{B,\diamond}
  \end{aligned}%
}{PowerSemigroupSingletonLayerReconstruction}{%
  \DeltaRow{Mengen}{A\dsep B\dsep a\dsep b}
  \DeltaRow{Funktionen}{\star\dsep\diamond\dsep\star_{\mathcal P}
    \dsep\diamond_{\mathcal P}\dsep F\dsep f\dsep\eta_A\dsep\eta_B}
}
\begin{proof}
Die Einschränkung von \(F\) auf \(\operatorname{Sing}(A)\) ist wegen der
Bildgleichheit eine Bijektion auf \(\operatorname{Sing}(B)\).  Daher definiert
\[
  f\coloneqq
  \eta_B^{-1}\circ
  \bigl(F\upharpoonright\operatorname{Sing}(A)\bigr)\circ\eta_A
  \colon A\bij B
\]
eine Bijektion mit \(F(\{a\})=\{f(a)\}\) für alle \(a\in A\).  Für
\(a,b\in A\) folgt
\[
\begin{aligned}
  \{f(a\star b)\}
  &=F(\{a\star b\})
   =F(\{a\}\star_{\mathcal P}\{b\})\\
  &=F(\{a\})\diamond_{\mathcal P}F(\{b\})
   =\{f(a)\diamond f(b)\}.
\end{aligned}
\]
Somit gilt \(f(a\star b)=f(a)\diamond f(b)\), und \(f\) ist ein
Halbgruppenisomorphismus.
\end{proof}

Die Erhaltung der Einermengenschicht ist hinreichend, aber nicht notwendig
für die Isomorphie der Ausgangshalbgruppen.  In der Potenzhalbgruppe der
zweielementigen Nullhalbgruppe ist jedes Produkt gleich \(\{0\}\).  Deshalb
ist
\[
  \{1\}\longleftrightarrow\{0,1\},
  \qquad \{0\}\longmapsto\{0\},
\]
ein Automorphismus der Potenzhalbgruppe, obwohl er eine Einermenge mit einer
zweielementigen Menge vertauscht.  Ausgangs- und Zielhalbgruppe sind dennoch
identisch.  Bei einem abstrakten Isomorphietest dürfen die Blöcke \(P_r\)
daher nicht als zusätzliche Farben verwendet werden.

\section{Rekonstruierbare Grundträgergröße}

Nach
\FormulaRefAuto{FiniteNonemptyPowerSetReflectsCardinality}[thm]
gilt für endliche Mengen \(A\) und \(B\)
\[
  \PowNE{A}\EqCard\PowNE{B}
  \quad\Longrightarrow\quad
  \tCard(A)=\tCard(B).
\]
Dabei gilt
\(\tCard(\PowNE A)+1=2^{\tCard(A)}\), und die
Zweierpotenzfunktion ist injektiv.  Insbesondere bestimmt jeder Isomorphismus endlicher
Potenzhalbgruppen die Zahl der Grundelemente, obwohl die einzelnen
Mächtigkeitsklassen im Allgemeinen nicht invariant sind.

\section{Der Fall zweier Grundelemente}

\FormulaThmDeltaK[Potenzhalbgruppen bestimmen Halbgruppen der Ordnung zwei]{%
  \begin{aligned}[t]
    &\FiniteSemigroupAlg{A}{\star}\dsep
      \FiniteSemigroupAlg{B}{\diamond}\dsep
      \tCard(A)=2\dsep\tCard(B)=2\dsep{}\\[-2pt]
    &\SemigroupIso{F}
      {\PowNE{A},\star_{\mathcal P}}
      {\PowNE{B},\diamond_{\mathcal P}}
      \vdash
      \exists f\,\FiniteSemigroupIso{f}{A,\star}{B,\diamond}
  \end{aligned}%
}{FinitePowerSemigroupReconstructionOrderTwo}{%
  \DeltaRow{Mengen}{A\dsep B\dsep S_2\dsep X\dsep Y}
  \DeltaRow{Funktionen}{\star\dsep\diamond\dsep F\dsep f\dsep m}
}
\begin{proof}
Jede zweielementige Halbgruppe wird durch eine Nummerierung
\(e\colon S_2\bij A\) auf eine assoziative Operation auf
\(S_2=\{0,1\}\) transportiert.  Nach
\FormulaRefAuto{FiniteSemigroupRelabelOrbitsAreIsomorphismClasses}[thm]
entsprechen deren Umbenennungsorbits genau den Isomorphieklassen.
Wir schreiben eine Tafel auf \(S_2\) zeilenweise als
\[
  t=(t_{00},t_{01},t_{10},t_{11})\in\{0,1\}^{4}.
\]
Genau die folgenden acht Tupel sind assoziativ:
\[
\begin{array}{cccc}
0000&0001&0011&0101\\
0110&0111&1001&1111.
\end{array}
\]
Für die konstanten Operationen sind beide Klammerungen gleich dem konstanten
Wert; für die linke beziehungsweise rechte Projektion sind sie beide gleich
\(x\) beziehungsweise \(z\).  Ferner gilt
\[
  \min(\min(x,y),z)=\min(x,y,z)=\min(x,\min(y,z)),
\]
und entsprechend für das Maximum.  Das Tupel \(0110\) beschreibt
\(x\star y\equiv x+y\pmod2\), das Tupel \(1001\) beschreibt
\(x\star y\equiv x+y+1\pmod2\).  In beiden Fällen ergeben die zwei
Klammerungen denselben Rest modulo \(2\).
Für jedes der übrigen acht Tupel liefert das angegebene Tripel verschiedene
Werte der beiden Klammerungen:
\[
\begin{array}{c|c}
\text{Tafeltupel}&\text{Tripel }(x,y,z)\text{ mit }(xy)z\neq x(yz)\\
\hline
0010,\ 0100&(1,0,1)\\
1000,\ 1110&(0,0,1)\\
1010,\ 1011,\ 1100,\ 1101&(0,0,0).
\end{array}
\]
Damit sind alle \(2^4=16\) Tafeln erfasst.

Die nichttriviale Umbenennung \(0\leftrightarrow1\) zerlegt die acht
assoziativen Tafeln in genau die fünf Orbits
\[
\{0000,1111\},\quad
\{0001,0111\},\quad
\{0011\},\quad
\{0101\},\quad
\{0110,1001\}.
\]
Sie entsprechen der konstanten Halbgruppe, dem Zweierhalbverband, der
Linksnullhalbgruppe, der Rechtsnullhalbgruppe und der Zweiergruppe.

Alle zugehörigen Potenzhalbgruppen besitzen drei Elemente.  Die folgenden
Eigenschaften hängen nur von ihrer abstrakten Multiplikation ab:
\begin{center}
\small
\begin{tabularx}{\linewidth}{@{}>{\raggedright\arraybackslash}Xc
  >{\raggedright\arraybackslash}X@{}}
\toprule
\textbf{Grundtyp und Regel}&
\textbf{Idempotente in \(\PowNE{S}\)}&
\textbf{Kennzeichen der Potenzhalbgruppe}\\
\midrule
Nullhalbgruppe (\(xy=0\))&1&konstant und kommutativ\\
Linksnullhalbgruppe (\(xy=x\))&3&\(XY=X\), nicht kommutativ\\
Rechtsnullhalbgruppe (\(xy=y\))&3&\(XY=Y\), nicht kommutativ\\
Zweierhalbverband (\(xy=\max\{x,y\}\))&3&kommutativ\\
Zweiergruppe (\(C_2\))&2&kommutativ\\
\bottomrule
\end{tabularx}
\end{center}
Die Anzahl der Idempotenten trennt zunächst die konstante Halbgruppe und die
Zweiergruppe ab.  Von den drei übrigen Fällen ist nur der Halbverband
kommutativ.  Links- und Rechtsnullfall sind ebenfalls getrennt: Wäre
\(F\) ein Isomorphismus vom Linksnull- zum Rechtsnullfall, so gälte für
\(X\neq Y\)
\[
  F(X)=F(XY)=F(X)F(Y)=F(Y),
\]
im Widerspruch zur Injektivität.  Isomorphe Potenzhalbgruppen gehören daher
zum selben der fünf Grundtypen, und die Ausgangshalbgruppen sind isomorph.
\end{proof}

\section{Der endliche Suchweg}

Die bisherigen Algorithmen ergeben für jede fest gewählte Ordnung \(N\) ein
endliches Entscheidungsverfahren:
\begin{enumerate}
  \item Erzeuge durch assoziatives Backtracking genau einen kanonischen
    Vertreter jeder Halbgruppenisomorphieklasse auf \(S_N\).
  \item Konstruiere zu jedem Vertreter die Tafel seiner
    \(2^N-1\)-elementigen Potenzhalbgruppe.
  \item Berechne den kanonischen Tafelcode der Potenzhalbgruppe unter
    \emph{allen} Permutationen ihrer \(2^N-1\) Elemente.
  \item Gruppiere die Ausgangshalbgruppen nach diesem Potenztafelcode.
\end{enumerate}
Enthält eine solche Codeklasse zwei nichtisomorphe Ausgangshalbgruppen, so ist
ein endliches Gegenbeispiel gefunden.  Bleiben alle Codeklassen einelementig,
ist die Vermutung für genau diese Ordnung \(N\) bewiesen.

\begin{remark}[Ungefärbte Kanonisierung]
Beim dritten Schritt dürfen weder Teilmengengröße noch Einermengenstatus als
zusätzliche Farbe verwendet werden.  Eine solche Markierung würde nur
Isomorphismen zulassen, die bereits die gesuchte Einermengenschicht erhalten,
und damit die eigentliche Schwierigkeit aus dem Test entfernen.  Zulässig
sind nur Eigenschaften, die aus der abstrakten Multiplikation selbst folgen,
etwa Idempotenz oder die Multimengen der Eintragshäufigkeiten in Zeilen und
Spalten.
\end{remark}

Das Verfahren ist endlich, wächst aber sehr schnell.  Schon die
Potenzhalbgruppentafel hat
\[
  (2^N-1)^2
\]
Felder, und ein naiver kanonischer Code würde bis zu
\((2^N-1)!\) Umnummerierungen prüfen.  Das erklärt, weshalb eine
Computerprüfung kleiner Ordnungen wertvoll ist, aber keinen allgemeinen
Beweis ersetzt.

\section{Isomorphieproblem und endliche Potenzhalbgruppenvermutung}

In der Literatur heißt die folgende Frage überwiegend
\emph{Isomorphieproblem für Potenzhalbgruppen}.  Wir nennen ihre positive
endliche Fassung in diesem Band die \emph{endliche
Potenzhalbgruppenvermutung}.  Weil bei endlichem Grundträger jede Teilmenge
endlich ist, genügt hier durchgehend die bereits verwendete Konstruktion
\(\PowNE{A}\).

\begin{samepage}
\begin{center}
\fbox{\begin{minipage}{0.9\linewidth}
\textbf{Endliche Potenzhalbgruppenvermutung.}
Für endliche Halbgruppen \((A,\star)\) und \((B,\diamond)\) gilt
\[
\begin{aligned}
&\exists F\,
  \SemigroupIso{F}
    {\PowNE{A},\star_{\mathcal P}}
    {\PowNE{B},\diamond_{\mathcal P}}
\quad\Longrightarrow{}\\[-2pt]
&\hspace{35mm}
  \exists f\,\SemigroupIso{f}{A,\star}{B,\diamond}.
\end{aligned}
\]
\end{minipage}}
\end{center}
\end{samepage}

Die in Band~22 bewiesene umgekehrte Richtung gilt sogar ohne Endlichkeit:
Ein Isomorphismus \(f\colon A\to B\) induziert durch \(X\mapsto f[X]\) einen
Isomorphismus der Potenzhalbgruppen.  Nur die eingerahmte Richtung ist offen
und wird bewusst weder als Axiom noch als Theorem registriert.

Zwei Stärken der Rekonstruktionsfrage sind auseinanderzuhalten.

\smallskip
\noindent\textbf{Global bestimmt.}
Aus der Existenz irgendeines Isomorphismus der Potenzhalbgruppen folgt die
Existenz eines Isomorphismus der Grundhalbgruppen.

\smallskip
\noindent\textbf{Starke Isomorphieeigenschaft.}
Jeder konkrete Isomorphismus der Potenzhalbgruppen bildet die
Einermengenschicht auf die Einermengenschicht ab.

\smallskip
Die starke Eigenschaft impliziert die globale Bestimmtheit, aber die
Nullhalbgruppe zeigt, dass die Umkehrung nicht als Beweisstrategie vorausgesetzt
werden darf.  Eine \emph{einseitige Starrheitsaussage} verstärkt die globale
Bestimmtheit in einer anderen Richtung: Nur eine Grundhalbgruppe wird in
einer speziellen Klasse vorausgesetzt, die andere darf zunächst beliebig
sein.  Mit der starken Isomorphieeigenschaft ist dies im Allgemeinen nicht
vergleichbar: Die einseitige Aussage kontrolliert die Zielklasse, die starke
Eigenschaft dagegen jeden konkreten Isomorphismus zwischen
Potenzhalbgruppen.

Ohne Endlichkeitsvoraussetzung ist die globale Aussage falsch.  Das
vollständig bewiesene unendliche Gegenbeispiel steht in
\FormulaRefAuto{MogiljanskajaPairCounterexample}[thm].  Es entscheidet den
endlichen Fall nicht.

\subsection{Endliche links-, rechts- und beidseitig kürzbare Fälle}

Für die in den Bänden~27--29 eingeführten Kürzbarkeitsklassen lässt sich die
endliche Vermutung entscheiden.  Der elementare Teil der Argumentation wird
zunächst vollständig bewiesen; beim anschließenden Rekonstruktionsschritt
wird jeweils der genau bezeichnete Literatursatz verwendet.

\paragraph{Hilfssatz: Endliche kürzbare Halbgruppen.}
Jede nichtleere endliche Halbgruppe, die links- und rechtskürzbar ist, ist
eine Gruppe.

\emph{Beweis.}
Sei \((A,\star)\) eine solche Halbgruppe, und wähle \(a_0\in A\).  Die
Translationen
\[
  L_{a_0}\colon A\to A,\quad x\longmapsto a_0\star x,
  \qquad
  R_{a_0}\colon A\to A,\quad x\longmapsto x\star a_0
\]
sind wegen der Links- beziehungsweise Rechtskürzbarkeit injektiv und wegen
der Endlichkeit von \(A\) bijektiv.  Daher gibt es \(e_\ell,e_r\in A\) mit
\[
  a_0\star e_\ell=a_0,
  \qquad
  e_r\star a_0=a_0.
\]
Für jedes \(x\in A\) gilt wegen der Assoziativität
\[
  a_0\star(e_\ell\star x)
  =(a_0\star e_\ell)\star x
  =a_0\star x.
\]
Linkskürzung liefert \(e_\ell\star x=x\); also ist \(e_\ell\) ein linkes
neutrales Element auf ganz \(A\).  Dual folgt aus
\[
  (x\star e_r)\star a_0
  =x\star(e_r\star a_0)
  =x\star a_0
\]
durch Rechtskürzung \(x\star e_r=x\); somit ist \(e_r\) ein rechtes
neutrales Element.  Folglich
\[
  e_\ell=e_\ell\star e_r=e_r.
\]
Schreibe \(e\coloneqq e_\ell=e_r\).  Für beliebiges \(x\in A\) sind auch
\(L_x\) und \(R_x\) bijektiv; daher gibt es \(b,c\in A\) mit \(x\star b=e\) und
\(c\star x=e\).  Dann ist
\[
  b=e\star b=(c\star x)\star b
    =c\star(x\star b)=c\star e=c.
\]
Damit besitzt jedes \(x\in A\) ein beidseitiges Inverses, und
\((A,\star)\) ist eine Gruppe. \(\square\)

\paragraph{Struktursatz: Endliche einseitige Kürzbarkeit.}
Jede nichtleere endliche linkskürzbare Halbgruppe ist vollständig regulär,
also eine Vereinigung von Untergruppen.  Dasselbe gilt dual für jede
nichtleere endliche rechtskürzbare Halbgruppe.

\emph{Beweis.}
Sei zunächst \((A,\star)\) endlich und linkskürzbar, und sei \(a\in A\).
Die von \(a\) erzeugte Unterhalbgruppe
\[
  \langle a\rangle_{A,\star}
  =\{a^n\mid n\geq 1\}
\]
ist nichtleer, endlich und linkskürzbar.
Außerdem kommutieren ihre Elemente, denn für positive natürliche Zahlen
\(m,n\) gilt
\[
  a^m\star a^n=a^{m+n}=a^{n+m}=a^n\star a^m.
\]
In einer kommutativen Halbgruppe folgt aus der Linkskürzbarkeit auch die
Rechtskürzbarkeit: Aus \(x\star z=y\star z\) wird durch Vertauschen der
Faktoren \(z\star x=z\star y\), und Linkskürzung ergibt \(x=y\).
Der Hilfssatz zeigt daher, dass \(\langle a\rangle_{A,\star}\) eine Gruppe
ist.  Insbesondere liegt \(a\) in einer Untergruppe von \(A\).  Da dies für
jedes \(a\in A\) gilt, ist \(A\) vollständig regulär.

Ist \((A,\star)\) stattdessen rechtskürzbar, so ist
\(\langle a\rangle_{A,\star}\) rechtskürzbar und wiederum kommutativ.  Aus
\(z\star x=z\star y\) folgt dann
\(x\star z=y\star z\); Rechtskürzung liefert \(x=y\).  Die erzeugte
Unterhalbgruppe ist somit auch linkskürzbar, nach dem Hilfssatz eine Gruppe
und enthält \(a\).  Also ist auch jede nichtleere endliche rechtskürzbare
Halbgruppe vollständig regulär. \(\square\)

\paragraph{Folgerung für Potenzhalbgruppen.}
Seien \((A,\star)\) und \((B,\diamond)\) endliche Halbgruppen mit
\[
  \PowNE{A}\cong\PowNE{B}.
\]
Dann folgt \((A,\star)\cong(B,\diamond)\) in jedem der folgenden drei
Fälle:
\begin{enumerate}
  \item Beide Halbgruppen sind linkskürzbar.
  \item Beide Halbgruppen sind rechtskürzbar.
  \item Mindestens eine der beiden Halbgruppen ist links- und
    rechtskürzbar; für die andere ist keine Kürzbarkeit vorausgesetzt.
\end{enumerate}

Sind die Träger nichtleer, so sind in den ersten beiden Fällen beide
Halbgruppen nach dem Struktursatz vollständig regulär.  Der Satz über die
globale Bestimmtheit vollständig regulärer Halbgruppen von
\href{https://doi.org/10.1080/00927872.2026.2659846}%
  {Yu--Zhao (2026), Satz~5.4}
liefert daher ihre Isomorphie.  Die zugängliche Fassung dieser Arbeit arbeitet
allgemein unter der verallgemeinerten Kontinuumshypothese.  Der betreffende
Kardinalitätsschritt ist hier jedoch ohne diese Zusatzannahme gültig: Für die
im Beweis auftretenden Klassen
\(C\subseteq A\) und \(D\subseteq B\) folgt aus einer Bijektion
\(\PowNE{C}\to\PowNE{D}\) wegen der Endlichkeit
\[
  2^{|C|}-1=2^{|D|}-1,
  \qquad\text{also}\qquad
  |C|=|D|.
\]
Damit spezialisiert sich der Rekonstruktionsbeweis auf endliche Träger ohne
eine mengen-theoretische Zusatzannahme.  Im dritten Fall ist die kürzbare
Halbgruppe nach dem Hilfssatz eine Gruppe.  Der in Band~33 vollständig
bewiesene einseitige Gruppenstarrheitssatz von
\href{https://doi.org/10.48550/arXiv.2606.01917}%
  {Liu--Tringali (2026)}
besagt gerade, dass aus der Isomorphie ihrer großen Potenzhalbgruppe mit der
Potenzhalbgruppe einer beliebigen Halbgruppe bereits die Isomorphie der
Grundhalbgruppen folgt.

Der im lokalen Formalismus mögliche leere Träger verursacht keine Ausnahme:
Es gilt \(\PowNE{A}=\varnothing\) genau dann, wenn \(A=\varnothing\).  Ein
Isomorphismus der Potenzhalbgruppen erzwingt daher, dass entweder beide
Träger leer oder beide nichtleer sind; zwei leere Halbgruppen sind durch die
leere Abbildung isomorph.  Damit ist die Folgerung in allen drei Fällen
bewiesen.

Die Aussage ist eine globale Bestimmtheitsaussage, nicht die starke
Isomorphieeigenschaft.  Das ist bereits bei endlichen einseitig kürzbaren
Halbgruppen wesentlich: Für \(x\star y=y\) ist die Grundhalbgruppe
linkskürzbar und ihre Potenzhalbgruppe erfüllt \(X\star_{\mathcal P}Y=Y\).
Somit ist jede Bijektion ihrer Potenzhalbgruppe ein Automorphismus; bei
mindestens zwei Grundelementen kann ein solcher Automorphismus eine
Einermenge mit einer Nichteinermenge vertauschen.  Für rechtskürzbare
Halbgruppen liefert \(x\star y=x\) das duale Beispiel.

\subsection{Bekannte Resultate und Status}

\begin{remark}
Die Vermutung ist mit Stand vom 28.~August~2026 für beliebige endliche
Halbgruppen offen; siehe
\href{https://doi.org/10.1007/978-3-031-75326-8_20}%
  {Tringali, \emph{On the Isomorphism Problem for Power Semigroups}, 2025}
und
\href{https://doi.org/10.48550/arXiv.2606.01917}%
  {Liu--Tringali, \emph{Power Semigroups and Two Rigidity Theorems for
  Groups}, 2026}.
Tamuras früherer Beitrag
\href{https://doi.org/10.1007/978-94-009-3839-7_22}%
  {\emph{On the Recent Results in the Study of Power Semigroups}, 1987}
berichtet die positive Aussage ausdrücklich ohne Beweis und wird hier daher
nicht als Lösung verwendet.

Wichtige bekannte Teilresultate sind:
\begin{itemize}
  \item Ohne Endlichkeitsvoraussetzung ist die allgemeine Aussage falsch.
    Das in Band~22 vollständig ausgeführte Gegenbeispiel geht zurück auf
    \href{https://doi.org/10.1007/BF02389140}%
      {Mogiljanskaja, \emph{Non-isomorphic Semigroups with Isomorphic
      Semigroups of Subsets}, 1973}.
  \item Sind beide Halbgruppen kürzbar und ist mindestens eine von ihnen
    kommutativ, so gilt sogar die starke Isomorphieeigenschaft
    \href{https://doi.org/10.1007/978-3-031-75326-8_20}%
      {(Tringali, 2025)}.  Das Resultat reicht über den endlichen Rahmen
      hinaus.
  \item Sind beide endlichen Halbgruppen idempotent, so folgt ihre
    Isomorphie aus
    \FormulaRefAuto{FiniteIdempotentSemigroupsGloballyDetermined}[thm].
    Band~24 enthält den Rekonstruktionsbeweis und erläutert auch, weshalb
    die in der Literatur unbeschränkt formulierte unendliche Fassung nicht
    ohne einen zusätzlichen kardinalarithmetischen Schluss übernommen
    werden darf.
  \item Sind beide Halbgruppen vollständig regulär, also Vereinigungen von
    Gruppen, so ist diese Fachklasse nach der zugänglichen Fassung von
    Yu--Zhao unter der verallgemeinerten Kontinuumshypothese global bestimmt
    \href{https://doi.org/10.1080/00927872.2026.2659846}%
      {(Yu--Zhao, 2026)}.  Für endliche Träger ist die in der vorangehenden
      Untersektion ausgeführte Anwendung davon unabhängig.
  \item Ist eine Grundhalbgruppe eine Gruppe, darf die andere zunächst eine
    beliebige Halbgruppe sein; aus isomorphen Potenzhalbgruppen folgt
    dennoch die Isomorphie der Grundhalbgruppen
    \href{https://doi.org/10.48550/arXiv.2606.01917}%
      {(Liu--Tringali, 2026)}.  Der vollständige Beweis steht in Band~33 im
      Abschnitt über die Rekonstruktion aus der großen Potenzhalbgruppe.
\end{itemize}
Diese positiven Inseln erklären mögliche Rekonstruktionsstrategien, entscheiden
die Vermutung für beliebige endliche Halbgruppen aber nicht.
\end{remark}


%============================================================
%============================================================
\chapter{Rechnerischer Anhang}
%============================================================

Die folgenden Abschnitte betreffen die vollständige Erzeugung und
Identifikation kleiner Multiplikationstafeln.  Sie sind rechnerische
Anwendungen der zuvor bewiesenen Struktur- und Orbit-Aussagen und werden
für deren Beweise nicht benötigt.

\section{Rekursive Erzeugung von Halbgruppentafeln}
%============================================================

\subsection{Was die Potenzrekursion tatsächlich erzeugt}

Zwei Verfahren sind voneinander zu unterscheiden.  Die
\emph{Potenzrekursion} verfolgt die von einem Element erzeugte
Unterhalbgruppe.  Ein \emph{Backtracking-Verfahren} versucht dagegen, eine
partielle Multiplikationstafel zu einer vollständigen assoziativen Tafel
fortzusetzen.  Das erste Verfahren erklärt die Struktur des folgenden
Leitbeispiels; erst das zweite Verfahren erfasst beliebige endliche
Halbgruppen.

Sei \(N\in\mathbb N\) mit \(N\geq1\) und
\[
  S_N\coloneqq\NatSeg{N}=\{0,\ldots,N-1\}
\]
und sei zunächst
\[
  m\colon S_N\times S_N\to S_N
\]
bereits die assoziative Multiplikationstafel einer Halbgruppe.  Ausgehend von
dem Element \(0\) definieren wir
\[
  x_1\coloneqq0,
  \qquad
  x_{r+1}\coloneqq m(x_r,0)
  \quad(r\geq1).
\]
Dann ist \(x_r\) die \(r\)-te Potenz von \(0\).  Das Additionsgesetz für
positive Halbgruppenpotenzen aus Band~22 liefert unmittelbar
\[
  \boxed{m(x_r,x_s)=x_{r+s}}
  \qquad(r,s\geq1).
\]
Insbesondere gilt
\[
  m(x_r,0)=m(0,x_r)=x_{r+1}.
\]
Nach der Definition und der Potenzbeschreibung der monogenen
Unterhalbgruppe aus Band~22 gilt
\[
  \langle a\rangle_{S_N,m}
  =P_{S_N,m}(a)
  =\{a^r\mid r\in\mathbb N,\ r\neq0\}.
\]
Dabei werden ausdrücklich
\FormulaRefAuto{SemigroupMonogenicSubsemigroupDef}[def] und
\FormulaRefAuto{SemigroupMonogenicPowerDescription}[thm] verwendet.  Eine
Halbgruppe heißt \emph{monogen}, wenn ihr Träger gleich
\(\langle a\rangle_{S_N,m}\) für ein Element \(a\) ist.  Im Folgenden wird
also \(\langle0\rangle_{S_N,m}\) untersucht.  Der Großbuchstabe \(N\) bezeichnet dabei
stets die Trägergröße; er entspricht dem zuvor verwendeten Parameter \(n\)
und damit der Mächtigkeit des Trägers \(\NatSeg{n}\).

In der von \(0\) erzeugten Unterhalbgruppe liegen die Tafeleinträge daher
auf »Diagonalen mit gleicher Indexsumme«:
\[
\begin{array}{c|ccccc}
\star &x_1&x_2&x_3&x_4&\cdots\\\hline
x_1&x_2&x_3&x_4&x_5&\cdots\\
x_2&x_3&x_4&x_5&x_6&\cdots\\
x_3&x_4&x_5&x_6&x_7&\cdots\\
x_4&x_5&x_6&x_7&x_8&\cdots\\
\vdots&\vdots&\vdots&\vdots&\vdots&\ddots
\end{array}
\]

Eine wichtige Konsistenzbedingung wird dadurch sofort sichtbar.  Sobald
zwei Potenzen übereinstimmen, müssen auch alle ihre Fortsetzungen
übereinstimmen:
\[
  x_r=x_s
  \quad\Longrightarrow\quad
  x_{r+t}=x_{s+t}
  \qquad(t\geq1).
\]
Der Index--Perioden-Satz dieses Bandes liefert wegen der Endlichkeit von
\(S_N\) eindeutig bestimmte Zahlen \(\mu,\lambda\geq1\), so dass
\[
  x_1,\ldots,x_{\mu+\lambda-1}
  \quad\text{paarweise verschieden sind}
  \qquad\text{und}\qquad
  x_{\mu+\lambda}=x_\mu.
\]
Die Zahl \(\mu\) ist der Beginn des periodischen Teils und \(\lambda\) dessen
Periode.  Die bereits bewiesene periodische Fortsetzung zeigt, dass die Folge
aus einem Vorlauf und einem anschließenden Zyklus besteht:
\[
  \underbrace{x_1,\ldots,x_{\mu-1}}_{\text{Vorlauf}},
  \quad
  \underbrace{x_\mu,\ldots,x_{\mu+\lambda-1}}_{\text{Zyklus}},
  \quad
  x_{\mu+\lambda}=x_\mu.
\]
Mit
\[
  \rho_{\mu,\lambda}(t)
  \coloneqq
  \begin{cases}
    t,&t<\mu,\\
    \mu+\bigl((t-\mu)\bmod\lambda\bigr),&t\geq\mu
  \end{cases}
\]
lässt sich jede Potenz auf ihren eindeutigen Vertreter reduzieren:
\[
  \boxed{x_t=x_{\rho_{\mu,\lambda}(t)}}\qquad(t\geq1).
\]
Für \(t<\mu\) ist dies die Identität.  Sei also \(t\geq\mu\).  Nach dem
Divisionsalgorithmus existieren eindeutig \(q,r\in\mathbb N\) mit
\[
  t-\mu=q\lambda+r,\qquad r<\lambda.
\]
Die Verschiebung um \(q\) Perioden ergibt
\[
  x_t=x_{\mu+r}
     =x_{\rho_{\mu,\lambda}(t)}.
\]
Damit gilt auf der monogenen Unterhalbgruppe
\[
  \boxed{
    m(x_r,x_s)=x_{\rho_{\mu,\lambda}(r+s)}
  }
\]
Index und Periode bestimmen somit die gesamte Tafel von
\[
  \langle0\rangle_{S_N,m}=\{x_r\mid r\geq1\}.
\]

Die Aussage besitzt auch eine Konstruktionsrichtung.  Wählt man beliebige
\(\mu,\lambda\geq1\), nimmt
\[
  x_1,\ldots,x_{\mu+\lambda-1}
\]
als verschiedene Symbole und \emph{definiert} ihre Multiplikation durch die
Reduktionsformel, so entsteht auf einem Träger mit
\(\mu+\lambda-1\) Elementen eine assoziative monogene Halbgruppe.

Liegt genau eine der positiven natürlichen Zahlen \(u,v\) unter \(\mu\), so
liegt auch genau einer der beiden Reduktionswerte unter \(\mu\).  Liegen beide
Zahlen unter \(\mu\), so gilt
\(\rho_{\mu,\lambda}(u)=\rho_{\mu,\lambda}(v)\) genau für \(u=v\).
Liegen beide mindestens bei \(\mu\), so ist die Gleichheit der
Reduktionswerte äquivalent zur Gleichheit der Divisionsreste von
\(u-\mu\) und \(v-\mu\) modulo \(\lambda\).  Folglich gilt
\[
  \rho_{\mu,\lambda}(u)=\rho_{\mu,\lambda}(v)
  \quad\Longleftrightarrow\quad
  \left(
    \begin{array}{l}
      u=v<\mu,\ \text{oder}\\
      u,v\geq\mu\ \text{und}\ u\equiv v\pmod\lambda
    \end{array}
  \right).
\]
Die rechts beschriebene Relation ist mit der Addition positiver natürlicher
Zahlen verträglich: Gleichheit bleibt nach Addition erhalten, und Zahlen, die
mindestens \(\mu\) und modulo \(\lambda\) kongruent sind, behalten beide
Eigenschaften.  Da \(u\) und \(\rho_{\mu,\lambda}(u)\) stets in derselben
Klasse liegen, folgt für \(u,t\geq1\)
\[
  \rho_{\mu,\lambda}\bigl(\rho_{\mu,\lambda}(u)+t\bigr)
  =
  \rho_{\mu,\lambda}(u+t).
\]
Wegen der Kommutativität der Addition gilt ebenso
\[
  \rho_{\mu,\lambda}\bigl(t+\rho_{\mu,\lambda}(u)\bigr)
  =
  \rho_{\mu,\lambda}(t+u).
\]
Insbesondere gilt
\[
  \rho_{\mu,\lambda}\!
    \bigl(\rho_{\mu,\lambda}(r+s)+t\bigr)
  =
  \rho_{\mu,\lambda}(r+s+t)
  =
  \rho_{\mu,\lambda}\!
    \bigl(r+\rho_{\mu,\lambda}(s+t)\bigr).
\]
In der folgenden kurzen Rechnung lassen wir das Operationszeichen weg und
schreiben \(x_rx_s\) für \(m(x_r,x_s)\).  Für \(r,s,t\geq1\) gilt damit:
\[
\begin{aligned}
  (x_r x_s)x_t
  &=x_{\rho(\rho(r+s)+t)}
   =x_{\rho(r+s+t)}\\
  &=x_{\rho(r+\rho(s+t))}
   =x_r(x_sx_t),
\end{aligned}
\]
wobei hier zur Entlastung \(\rho=\rho_{\mu,\lambda}\) geschrieben wurde.
Damit ist die Potenzrekursion ein vollständiger Erzeugungsalgorithmus für
monogene endliche Halbgruppen bis auf Isomorphie.  Soll die Tafel auf dem
fest nummerierten Träger \(S_N\) liegen, muss
\(N=\mu+\lambda-1\) gelten; zusätzlich wählt man eine Bijektion von
\(\{x_1,\ldots,x_N\}\) nach \(S_N\).  Alle daraus entstehenden
Umnummerierungen liefert der spätere Permutationsalgorithmus.  Für eine feste
Ordnung \(N\) kann man daher
\[
  \mu=1,\ldots,N,
  \qquad
  \lambda=N+1-\mu
\]
der Reihe nach einsetzen.  So entstehen genau die \(N\) monogenen
Isomorphietypen dieser Ordnung.  Für \(\mu>1\) ist \(x_1\) der einzige
Erzeuger: Für \(k\geq2\) und \(q\geq1\) gilt
\[
  x_k^q=x_{\rho_{\mu,\lambda}(kq)}\neq x_1.
\]
Tatsächlich ist im Fall \(kq<\mu\) der Reduktionswert \(kq\geq2\), während
er im Fall \(kq\geq\mu\) mindestens \(\mu>1\) ist.

Für \(\mu=1\) ist \(x_N\) neutral, denn
\[
  x_Nx_s=x_{\rho_{1,N}(N+s)}=x_s
  =x_{\rho_{1,N}(s+N)}=x_sx_N.
\]
Erzeugt ein Element \(x_k\) die ganze Halbgruppe, so existiert ein
\(q\geq1\) mit \(x_k^q=x_N\). Wähle das kleinste solche \(q\). Dann ist
\(x_k^{q+1}=x_k\). Die Potenzen
\[
  x_k,x_k^2,\ldots,x_k^q
\]
sind paarweise verschieden. Denn wären \(x_k^r=x_k^s\) und
\(1\leq r<s\leq q\), so widerspräche im Fall
\(s=q\) bereits \(x_k^r=x_k^q=x_N\) der Minimalität von \(q\). Im Fall
\(s<q\) ist \(q-s\geq1\); durch Multiplikation mit \(x_k^{q-s}\) und das
Potenzadditionsgesetz folgte dann
\[
  x_k^{r+q-s}=x_k^q=x_N,
\]
obwohl \(1\leq r+q-s<q\); dies widerspräche der Minimalität von \(q\).
Wegen \(x_k^{q+1}=x_k\) wiederholen sich alle späteren Potenzen mit Periode
\(q\). Da \(x_k\) die ganze \(N\)-elementige Halbgruppe erzeugt, bestehen
seine Potenzen damit aus genau den obigen \(q\) Elementen, also ist \(q=N\).
Folglich besitzt die Potenzfolge das minimale Index--Perioden-Paar
\((1,N)\).

Ein Isomorphismus bildet Erzeuger auf Erzeuger ab und erhält deren
Potenzgleichheiten.  Für \(\mu>1\) hat der einzige Erzeuger daher dasselbe
Index--Perioden-Paar; für \(\mu=1\) ist dieses Paar stets \((1,N)\).
Verschiedene Paare liefern folglich nichtisomorphe Halbgruppen.

\begin{remark}[Analyse und Konstruktion]
Aus der Vorgabe \(m(0,0)=k\) folgt zunächst \(x_2=k\).  Ist damit bereits
eine Wiederholung eingetreten, etwa \(k=0\), so ist die weitere Potenzfolge
erzwungen.  Nur solange noch keine Wiederholung vorliegt und der nächste Wert
nicht durch andere Tafeleinträge feststeht, entsteht bei der Konstruktion
eine echte Wahl.  Die Assoziativität überträgt jede solche Wahl auf weitere
Tafelfelder.  Die Formel \(m(x_r,x_s)=x_{r+s}\) ist also eine notwendige
Folgerung aus einer assoziativen Tafel und zugleich eine Konstruktionsregel
für den monogenen Teil, aber für sich allein noch kein Algorithmus zur
Erzeugung jeder Halbgruppe auf \(S_N\).
\end{remark}


%============================================================
\subsection{Leitbeispiel: Von der Potenzfolge zur ganzen Tafel}
%============================================================

Auf \(S_3=\{0,1,2\}\) beginnen wir mit
\[
  m(0,0)=2.
\]
Die ersten Schritte werden in der folgenden Tabelle getrennt als
Folgerungen und als Wahlen ausgewiesen:
\[
\begin{array}{c|c|l}
r&x_r&\text{Bedeutung}\\\hline
1&0&\text{Ausgangselement}\\
2&2&m(0,0)=2\\
3&1&m(2,0)=m(0,2)=1\quad\text{(gewählter gemeinsamer Wert)}\\
4&0&m(1,0)=m(0,1)=0\quad\text{(gewählter gemeinsamer Wert)}
\end{array}
\]
Da \(x_4=x_1\) gilt, erzwingt die Konsistenzbedingung
\[
  x_{r+3}=x_r
  \qquad(r\geq1).
\]
Hier sind also \(\mu=1\) und \(\lambda=3\).  Außerdem treten \(0,2,1\)
bereits als \(x_1,x_2,x_3\) auf.  Deshalb ist
\(\langle0\rangle_{S_3,m}=S_3\), und die Potenzrekursion bestimmt in diesem
besonderen Fall tatsächlich die ganze Tafel.

Die Wahl \(x_4=0\) ist jedoch nur einer von drei möglichen
Wiederholungspunkten.  Nachdem \(x_1=0,x_2=2,x_3=1\) feststehen, besitzt der
verbleibende Potenzrekursionsbaum für monogene Fortsetzungen auf \(S_3\) nur
noch drei Äste:
\[
\begin{array}{c|c|c}
\text{Wiederholung}&(\mu,\lambda)&\text{erzwungene Fortsetzung}\\\hline
x_4=x_1=0&(1,3)&0,2,1,0,2,1,\ldots\\
x_4=x_2=2&(2,2)&0,2,1,2,1,2,\ldots\\
x_4=x_3=1&(3,1)&0,2,1,1,1,1,\ldots
\end{array}
\]
Nach dieser einen Entscheidung ist jeweils jede weitere Potenz erzwungen.
Der erste Ast ist das Leitbeispiel; der dritte Ast wird im folgenden
Kontrastbeispiel vollständig ausgerechnet.

Zunächst ordnen wir Zeilen und Spalten in der Reihenfolge der Potenzfolge.
So ist die Herkunft jedes Eintrags unmittelbar sichtbar:
\[
\boxed{
\begin{array}{c|ccc}
\star&x_1=0&x_2=2&x_3=1\\\hline
x_1=0&x_2=2&x_3=1&x_4=x_1=0\\
x_2=2&x_3=1&x_4=x_1=0&x_5=x_2=2\\
x_3=1&x_4=x_1=0&x_5=x_2=2&x_6=x_3=1
\end{array}}
\]
Ordnet man anschließend Zeilen und Spalten numerisch als \(0,1,2\), erhält
man
\[
\begin{array}{c|ccc}
\star&0&1&2\\\hline
0&2&0&1\\
1&0&1&2\\
2&1&2&0
\end{array}.
\]
Das Element \(1\) ist neutral, und jedes Element besitzt bezüglich \(1\) ein
beidseitiges Inverses.  In der erst in Band~33 formal eingeführten
Gruppenterminologie ist die entstandene Halbgruppe die zyklische Gruppe mit
drei Elementen, nur in der Reihenfolge \(0,1,2\) ungewöhnlich nummeriert.


%============================================================
\subsection{Zwei Kontrastbeispiele}
%============================================================

Eine Potenzfolge muss nicht von Anfang an periodisch sein.  Wählen wir drei
verschiedene Symbole \(a,b,c\) und
\[
  x_1=a,\qquad x_2=b,\qquad x_3=c,\qquad x_4=c,
\]
so sind \(\mu=3\) und \(\lambda=1\).  Die Reduktionsformel ergibt die
assoziative Tafel
\[
\begin{array}{c|ccc}
\star&a&b&c\\\hline
a&b&c&c\\
b&c&c&c\\
c&c&c&c
\end{array}.
\]
Hier stehen \(a\) und \(b\) im Vorlauf; danach bleibt die Folge beim
Fixpunkt \(c\).  Dieses Beispiel rechnet den dritten Ast des vorherigen
Potenzrekursionsbaums vollständig aus.

Anders verhält es sich auf \(S_2=\{0,1\}\), wenn \(m(0,0)=0\) gilt.  Dann
ist \(x_r=0\) für jedes \(r\), so dass die Rekursion nichts über die Felder
mit dem Element \(1\) aussagt.  Beispielsweise sind beide Tafeln
\[
\begin{array}{c|cc}
\star&0&1\\\hline
0&0&0\\
1&0&0
\end{array}
\qquad\text{und}\qquad
\begin{array}{c|cc}
\star&0&1\\\hline
0&0&1\\
1&1&1
\end{array}
\]
assoziativ und besitzen dieselbe von \(0\) ausgehende Potenzfolge.  Links
steht die Nullhalbgruppe, rechts der zweielementige Halbverband mit der
Maximumsoperation, also eine kommutative idempotente Halbgruppe.  Dieses
Beispiel zeigt, warum zur Erzeugung allgemeiner Tafeln ein zusätzlicher
Fortsetzungsalgorithmus erforderlich ist.


%============================================================
\subsection{Assoziative Fortsetzung durch Backtracking}
%============================================================

Eine partielle Tafel \(p\) weist zunächst nur einigen Feldern
\((a,b)\in S_N^2\) einen Wert in \(S_N\) zu.  Sind für ein Tripel
\(a,b,c\in S_N\) bereits
\[
  p(a,b)=u
  \qquad\text{und}\qquad
  p(b,c)=v
\]
bekannt, so wird die Assoziativitätsbedingung zu der Feldgleichheit
\[
  \boxed{p(u,c)=p(a,v).}
\]
Sind beide Seiten bekannt und verschieden, ist die partielle Tafel
unmöglich.  Ist nur eine Seite bekannt, erhält die andere denselben Wert.
Sind beide Seiten noch unbekannt, werden die beiden Felder als gleich zu
behandelnde Felder vorgemerkt.  Technisch hält man die reflexive, symmetrische
und transitive Hülle aller vorgemerkten Feldgleichheiten, etwa mit einer
Union--Find-Struktur.  Neue Gleichheiten vereinigen ihre Klassen.  Jede Klasse
darf höchstens einen bereits gesetzten Wert tragen; beim Vereinigen zweier
Klassen mit verschiedenen Werten liegt ein Widerspruch vor.  Jede neue
Festlegung kann weitere solche Folgerungen auslösen.

Damit ergibt sich ein vollständiger rekursiver Algorithmus:
\begin{enumerate}
  \item Beginne mit einer leeren partiellen Tafel; vorgegebene Werte wie
    \(p(0,0)=k\) können sofort eingetragen werden.
  \item Prüfe alle Tripel und propagiere die erzwungenen Werte und
    Feldgleichheiten, bis keine neue Information entsteht.
  \item Widersprechen sich zwei erzwungene Werte, verwirf diesen Zweig und
    gehe zur letzten Wahl zurück.
  \item Sind alle Felder belegt, gib die Tafel aus.  Andernfalls wähle eine
    noch unbelegte Gleichheitsklasse von Feldern.
  \item Probiere für diese Klasse nacheinander jeden Wert aus \(S_N\) und
    rufe das Verfahren mit der erweiterten partiellen Tafel erneut auf.
\end{enumerate}
Die Potenzrekursion kann dabei als zusätzliche Propagationsregel verwendet
werden, sobald ein Diagonaleintrag gewählt wurde.  Ein Zustand besteht aus
einer Äquivalenzrelation auf den \(N^2\) Tafelfeldern und einer partiellen
Wertzuweisung auf ihren Klassen.  Seine Invariante lautet: Jede assoziative
vollständige Fortsetzung weist äquivalenten Feldern denselben Wert zu und
stimmt mit allen gesetzten Klassenwerten überein.

Sind \(p(a,b)=u\) und \(p(b,c)=v\) gesetzt, so muss jede assoziative
Fortsetzung \(p(u,c)=p(a,v)\) erfüllen.  Das Vereinigen dieser Feldklassen und
das Übertragen eines bereits gesetzten Wertes erhalten daher die Invariante;
zwei verschiedene Werte in derselben Klasse schließen jede Fortsetzung aus.
Beim Verzweigen über eine ungesetzte Klasse liegt jede Fortsetzung in genau
einem der \(N\) Wertzweige.  Jede Verzweigung vermindert die Zahl ungesetzter
Klassen, so dass die Tiefe höchstens \(N^2\) beträgt.  Am vollständigen Blatt
werden alle \(N^3\) Assoziativitätsgleichungen geprüft.  Der Algorithmus
terminiert somit und gibt genau alle assoziativen vollständigen Fortsetzungen
aus.  Die grobe Schranke für die Zahl der möglichen Blätter ist
\(N^{N^2}\).

\begin{example}[Der vollständige Suchbaum nach einer Vorgabe]
Auf \(S_2=\{0,1\}\) sei zunächst nur \(p(0,0)=1\) vorgegeben.  Das Tripel
\((0,0,0)\) fordert
\[
  p(1,0)=p(0,1).
\]
Bezeichne den noch unbekannten gemeinsamen Wert mit \(u\).  Nun gibt es genau
zwei Zweige.
\begin{itemize}
  \item Für \(u=0\) erzwingt das Tripel \((1,0,0)\) den Wert
    \(p(1,1)=1\).
  \item Für \(u=1\) erzwingt das Tripel \((0,0,1)\) ebenfalls
    \(p(1,1)=1\).
\end{itemize}
Damit besitzt die Vorgabe genau zwei assoziative Fortsetzungen:
\[
\begin{array}{c|cc}
\star&0&1\\\hline
0&1&0\\
1&0&1
\end{array}
\qquad\text{und}\qquad
\begin{array}{c|cc}
\star&0&1\\\hline
0&1&1\\
1&1&1
\end{array}.
\]
Links steht die zweielementige Gruppe mit neutralem Element \(1\), rechts
die konstante Halbgruppe mit konstantem Wert \(1\).  Das Beispiel enthält in
kleinstmöglicher Form eine Feldgleichheit, eine Verzweigung, die Propagation
in beiden Zweigen und zwei vollständige Blätter.
\end{example}


%============================================================
\section{Rekursive Erzeugung und Isomorphie}
%============================================================

In diesem Abschnitt werden drei verschiedene Suchziele verwendet.  Ihre
Zustände und Ausgaben sollten nicht miteinander verwechselt werden:
\begin{center}
\small
\begin{tabularx}{\textwidth}{@{}lXl@{}}
\toprule
\textbf{Ziel}&\textbf{rekursiver Zustand}&\textbf{Ausgabe}\\
\midrule
Beschriftungen einer festen Tafel&partielle Permutation&alle Tafeln \(m^\pi\)\\
Halbgruppen einer festen Ordnung&partielle assoziative Tafel&je ein kanonischer Vertreter\\
Isomorphismen \(m\to d\)&partielle injektive Abbildung&alle Isomorphismen \(f\)\\
\bottomrule
\end{tabularx}
\end{center}

\subsection{Alle isomorphen nummerierten Tafeln einer Halbgruppe}

Sei \(m\) eine feste Halbgruppentafel auf \(S_N\).  Alle Umbenennungen können
rekursiv erzeugt werden.  Man belegt dazu nacheinander
\(\pi(0),\pi(1),\ldots,\pi(N-1)\).  Im Schritt \(r\) wird für \(\pi(r)\)
jedes noch nicht verwendete Element von \(S_N\) ausprobiert.  Nach dem letzten
Schritt ist \(\pi\) eine Permutation, und man bildet die bereits definierte
umbenannte Tafel \(m^\pi\).  Nach dem Umbenennungskriterium sind \(m\) und
\(m^\pi\) isomorph, und jede
nummerierte Tafel, die zu \(m\) isomorph ist, entsteht auf diese Weise.

Im Leitbeispiel vertausche \(\pi\) die Zahlen \(0\) und \(1\) und lasse \(2\)
fest.  Die umbenannte, nun mit dem neutralen Element \(0\) beschriftete Tafel
lautet
\[
\begin{array}{c|ccc}
m^\pi&0&1&2\\\hline
0&0&1&2\\
1&1&2&0\\
2&2&0&1
\end{array}.
\]
Die beiden Tafeln sind verschieden beschriftet, aber isomorph.  Ein
\emph{Automorphismus} ist ein Isomorphismus einer Halbgruppe auf sich selbst.
Die Vertauschung von \(0\) und \(2\) bei festem \(1\) lässt dagegen die
ursprüngliche Tafel unverändert; sie ist ein Automorphismus.  Deshalb können
verschiedene Permutationen dieselbe umbenannte Tafel liefern.  Doppelte
Ausgaben vermeidet man, indem man die zeilenweise gelesenen Tafeltupel in
einer Menge speichert.  Die Anzahl verschiedener nummerierter Kopien ist
\[
  \frac{N!}{|\operatorname{Aut}(S_N,m)|}.
\]
Denn für \(\pi,\sigma\in\operatorname{Sym}(S_N)\) gilt nach dem
Wirkungsgesetz
\[
  m^\pi=m^\sigma
  \quad\Longleftrightarrow\quad
  \sigma^{-1}\circ\pi\in\operatorname{Aut}(S_N,m).
\]
Jede Faser der Abbildung \(\pi\mapsto m^\pi\) enthält daher genau
\(\lvert\operatorname{Aut}(S_N,m)\rvert\) Permutationen.  Die \(N!\)
Permutationen erhält man durch die aufeinanderfolgenden
Wahlmöglichkeiten \(N,N-1,\ldots,1\) für die Bilder der Trägerelemente.
Somit hat der Umbenennungsorbit die angegebene Größe.
Im Leitbeispiel gibt es \(3!=6\) Permutationen und genau zwei
Automorphismen.  Daher entstehen \(6/2=3\) verschiedene nummerierte Tafeln.


%============================================================
\subsection{Je ein Vertreter jeder Isomorphieklasse}
%============================================================

Soll nicht eine feste Halbgruppe umbenannt, sondern sollen \emph{alle}
Halbgruppen einer gegebenen Ordnung ohne isomorphe Mehrfachausgaben erzeugt
werden, verbindet man das Backtracking mit dem zuvor definierten kanonischen
Tafelcode:
\begin{enumerate}
  \item Erzeuge rekursiv alle assoziativen Tafeln auf \(S_N\).
  \item Berechne für jede vollständige Tafel \(m\) alle Tupel
    \(w_N(m^\pi)\).
  \item Gib \(m\) genau dann aus, wenn
    \[
      w_N(m)
      =
      \min_{\mathrm{lex}}
      \{w_N(m^\pi)\mid \pi\colon S_N\bij S_N\}.
    \]
\end{enumerate}
Jede Isomorphieklasse enthält eine solche lexikographisch kleinste
Nummerierung.  Umgekehrt können zwei verschiedene ausgegebene Tafeltupel
nicht isomorph sein, denn dann hätten sie denselben kanonischen Code.  Dieses
Verfahren liefert daher genau einen nummerierten Vertreter jeder
Isomorphieklasse.

Für einen zweielementigen Träger gibt es zunächst
\(2^{2^2}=16\) binäre Tafeln.  Acht davon sind assoziativ; der
Kanonizitätstest lässt genau fünf Isomorphieklassen mit den Tafeltupeln
\[
  (0,0,0,0),\quad (0,0,0,1),\quad (0,0,1,1),\quad
  (0,1,0,1),\quad (0,1,1,0)
\]
übrig.  Sie repräsentieren eine konstante Halbgruppe, den zweielementigen
Halbverband, die Linksnullhalbgruppe, die Rechtsnullhalbgruppe und die
zweielementige Gruppe.

Der Kanonizitätstest erst am vollständigen Blatt des Suchbaums ist nicht die
schnellste, aber die unmittelbar überprüfbare Variante.  Schnellere
Implementierungen verwenden eine kanonische Fortsetzung und verwerfen schon
partielle Tafeln, deren Nummerierung nicht kanonisch sein kann.  Dafür muss
jedoch zusätzlich bewiesen werden, dass die verwendete Regel in jeder
Isomorphieklasse mindestens einen Zweig bestehen lässt; der abschließende
Kanonizitätstest benötigt diese zusätzliche Theorie nicht.


%============================================================
\subsection{Isomorphismen zwischen zwei gegebenen Tafeln suchen}
%============================================================

Nach ihrer Nummerierung dürfen zwei zu vergleichende Tafeln \(m\) und \(d\)
auf demselben Träger \(S_N\) angenommen werden.  Auch alle Isomorphismen
zwischen ihnen lassen sich rekursiv bestimmen.  Der Zustand ist eine
partielle injektive Abbildung \(f\colon D\to S_N\) mit
\(D\subseteq S_N\).  Man wählt \(a\in S_N\setminus D\), probiert für \(f(a)\)
jedes noch freie Element von \(S_N\) und propagiert danach die
Homomorphiebedingung:
\[
  x,y\in D,\quad z=m(x,y)
  \quad\Longrightarrow\quad
  f(z)=d(f(x),f(y)).
\]
Ist \(z\) bereits abgebildet, muss sein Bild mit der rechten Seite
übereinstimmen.  Andernfalls wird \(z\) mit diesem Bild ergänzt, sofern das
Bild noch frei ist; bei einem Widerspruch wird zurückgegangen.  Nach jeder
Ergänzung wird diese Propagation bis zu einem Fixpunkt wiederholt.  Ist der
Definitionsbereich ganz \(S_N\), so ist die partielle Injektion eine
Bijektion und wegen der propagierten Produktgleichungen ein Isomorphismus.

Jeder Isomorphismus, der einen aktuellen Zustand erweitert, ordnet dem neu
gewählten Element \(a\) genau einen der noch freien Werte zu; der zugehörige
Zweig wird daher besucht.  Die Propagation fügt nur notwendige
Homomorphiegleichungen hinzu und entfernt keinen solchen Isomorphismus.  Bei
jeder Ergänzung wächst der endliche Definitionsbereich strikt.  Am
vollständigen Blatt sind alle Produktpaare geprüft.  Das Verfahren terminiert
also und gibt genau alle Isomorphismen zwischen \(m\) und \(d\) aus.

Zwischen der Ursprungstafel des Leitbeispiels und ihrer unmittelbar zuvor
angegebenen Umbenennung erzwingt die anfängliche Wahl \(f(0)=1\) wegen
\(0\star0=2\) zunächst \(f(2)=2\) und danach wegen \(2\star0=1\) den Wert
\(f(1)=0\).  So wächst der Isomorphismus schrittweise aus derselben Art von
Produktrekursion wie die ursprüngliche Potenzfolge.

\begin{remark}[Von der Potenzfolge zur Potenzhalbgruppentafel]
Die Potenzrekursion erzeugt \emph{nicht unmittelbar} die Potenzhalbgruppe.
Sie analysiert zunächst nur \(\langle0\rangle_{S_N,m}\) und liefert im monogenen Fall
die ganze Grundtafel \(m\).  Erst wenn diese Grundtafel vollständig feststeht,
werden Produkte nichtleerer Teilmengen gebildet und anschließend nummeriert:
\[
\boxed{
\begin{gathered}
  (x_r)\ \text{beziehungsweise}\ (\mu,\lambda)
  \ \longrightarrow\ \text{Grundtafel }m\\
  \longrightarrow\ X\star_{\mathcal P}Y
  \ \longrightarrow\ m_{\mathcal P}(p,q)
\end{gathered}
}
\]
Der zweite Pfeil bedeutet
\[
  X\star_{\mathcal P}Y
  =
  \{m(x,y)\mid x\in X,\ y\in Y\},
\]
der dritte lediglich die Umnummerierung durch \(\nu^{-1}\).  Im Leitbeispiel
liefert die Potenzfolge also zuerst die \(3\times3\)-Grundtafel; aus ihr
entsteht danach die \(7\times7\)-Tafel der sieben nichtleeren Teilmengen.
\end{remark}


%============================================================

\end{document}
